\chapter{The Radio Transceiver and the Software-Defined Radio}
\label{ch:sdr}

\begin{chapterintro}
Hold your phone in your hand and think about what its antenna is touching right now. FM stations
shouting tens of kilowatts from a hilltop, television transmitters, a dozen cellular carriers, your
neighbour's Wi-Fi, a microwave oven two rooms away, radar from the airport---and, somewhere in that
roar, the one base-station signal your phone actually wants, perhaps a hundred million times weaker
than the loudest of them. A radio receiver is a \emph{microphone for the electromagnetic spectrum}
that must pick out a single whisper in a stadium full of shouting, without being deafened by the
shouts and without adding hiss of its own.

The mathematics of the rest of this book assumes ideal mixers, perfect oscillators and linear
amplifiers. Real radios have none of these, and the engineer who understands their imperfections
designs systems that work on the first try. This chapter walks through a modern transceiver from the
antenna to the samples and back again. We start with the receiver architectures---the
superheterodyne, the direct-conversion receiver that now dominates integrated radios, the low-IF
compromise, image-reject mixers and direct RF sampling---and then take the RF front end apart:
filters and duplexers, the low-noise amplifier, the mixer. We develop the numbers every radio data
sheet quotes (sensitivity, noise figure, P1dB, IIP2, IIP3, spurious-free dynamic range, blocking) and
show how a receiver line-up is planned so that a $-100$\,dBm signal survives next to a $-25$\,dBm one.
We cover frequency generation from the quartz crystal to the fractional-$N$ synthesiser and Leeson's
model of phase noise, the I/Q impairments of direct conversion and the algorithms that remove them,
and the transmitter: DACs, power-amplifier classes, the efficiency--linearity trade-off, spectral
regrowth and ACLR, crest-factor reduction, digital predistortion, the Doherty amplifier and envelope
tracking.

The second half turns to the software-defined radio: its history from SPEAKeasy and Mitola to GNU
Radio and the \$20 dongle, its architecture as embodied in the Ettus USRP B200 on your bench, the
practical art of streaming samples between a radio and a computer, the measurements you can make with
it, and the platforms and industries that SDR has transformed. If you own a B200, keep it nearby:
almost every idea in this chapter can be seen on its screen.
\end{chapterintro}

%======================================================================
\section{Anatomy of a transceiver}
\label{sec:ch07:anatomy}
%======================================================================

\subsection{One skeleton, a billion radios}

Here is a surprising fact: a Bluetooth earbud, a smartphone, a Wi-Fi access point, a 64-antenna
5G base station and the USRP B200 on your bench all have the same skeleton. Their power levels differ
by a factor of a million and their prices by a factor of ten thousand, but draw their block diagrams
and you will draw the same picture (Figure~\ref{fig:ch07_transceiver}). Learn it once, and every radio
you ever open will look familiar.

Before we meet the blocks, look at the numbers that rule them (Figure~\ref{fig:ch07_by_numbers}).
Every one of them will reappear in this chapter, and by the end you should be able to explain each
one to a colleague over coffee.

\mdcfig{ch07_by_numbers}{The radio transceiver by the numbers. Each tile is explained somewhere in
this chapter: the thermal floor sets sensitivity, the dynamic range sets linearity, the filter count
is the scaling problem of the smartphone, the crystal sets frequency error, the converters set
bandwidth, the class-B limit sets battery life, and two inventions separated by 76 years---the
Doherty amplifier and the \$20 dongle---bracket the story.}

On the receive side, the skeleton is:
\begin{enumerate}
\item An antenna and a \emph{band filter} (a SAW, BAW, ceramic or cavity filter) that passes only the
band of interest and protects the receiver from strong out-of-band signals---including, in a
frequency-division duplex radio, its own transmitter.
\item A \term{low-noise amplifier} (LNA) that sets the noise figure of the whole chain
(Chapter~\ref{ch:noise}).
\item One or more \emph{mixers} driven by a \emph{local oscillator} (LO) that translate the signal to
an intermediate frequency or directly to baseband.
\item Analog filtering and variable gain (the \emph{baseband} or \emph{IF} section), under the control
of an automatic gain control loop.
\item A pair of ADCs, followed by digital down-conversion, decimation and channel filtering
(Chapter~\ref{ch:dsp}), and digital correction of the analog impairments.
\end{enumerate}
The transmit side mirrors it: digital predistortion and crest-factor reduction, a pair of DACs,
reconstruction filters, an I/Q up-converter, a driver, the power amplifier (PA), a harmonic filter,
and the antenna, shared with the receiver through a switch (time-division duplex) or a duplexer
(frequency-division duplex). A synthesiser, locked to a crystal reference, provides every LO and
every clock.

\begin{figure}[htbp]\centering
\begin{tikzpicture}[
  blk/.style={draw=navy,fill=navy!6,thick,minimum height=0.62cm,minimum width=1.05cm,align=center,font=\sffamily\scriptsize},
  dsp/.style={draw=accent,fill=accent!6,thick,minimum height=1.9cm,minimum width=1.5cm,align=center,font=\sffamily\scriptsize},
  syn/.style={draw=practc,fill=practc!8,thick,minimum height=0.62cm,minimum width=1.05cm,align=center,font=\sffamily\scriptsize},
  mix/.style={draw=navy,thick,circle,minimum size=0.5cm,inner sep=0pt,font=\sffamily\small},
  arr/.style={-{Stealth[length=1.6mm]},thick,navy},
  lo/.style={-{Stealth[length=1.4mm]},thick,practc}]
% antenna and duplexer
\node[draw=navy,thick,minimum height=4.4cm,minimum width=0.95cm,align=center,font=\sffamily\scriptsize] (dup) at (1.2,0) {dup-\\lexer\\or\\T/R\\switch};
\draw[thick,navy] (-0.1,0) -- (dup.west);
\draw[thick,navy] (-0.1,0) -- (-0.1,0.6) -- (-0.35,1.0) (-0.1,0.6) -- (0.15,1.0) (-0.35,1.0) -- (0.15,1.0);
\node[font=\sffamily\scriptsize] at (-0.1,1.25) {antenna};
% receive path
\node[blk] (lna) at (2.65,1.5) {LNA};
\node[mix] (mi) at (4.1,2.0) {$\times$};
\node[mix] (mq) at (4.1,1.0) {$\times$};
\node[blk] (fi) at (5.45,2.0) {LPF/VGA};
\node[blk] (fq) at (5.45,1.0) {LPF/VGA};
\node[blk] (adci) at (6.85,2.0) {ADC};
\node[blk] (adcq) at (6.85,1.0) {ADC};
\node[dsp] (drx) at (8.75,1.5) {DDC, decim.\\DC and I/Q\\correction\\AGC};
\draw[arr] (dup.east |- lna.west) -- (lna);
\draw[arr] (lna.east) -- (3.4,1.5) |- (mi.west);
\draw[arr] (3.4,1.5) |- (mq.west);
\draw[arr] (mi)--(fi); \draw[arr] (mq)--(fq); \draw[arr] (fi)--(adci); \draw[arr] (fq)--(adcq);
\draw[arr] (adci.east) -- (adci.east -| drx.west); \draw[arr] (adcq.east) -- (adcq.east -| drx.west);
% transmit path
\node[blk,fill=accent!6,draw=accent] (pa) at (2.65,-1.5) {PA};
\node[mix] (sum) at (3.4,-1.5) {$+$};
\node[mix] (ti) at (4.1,-1.0) {$\times$};
\node[mix] (tq) at (4.1,-2.0) {$\times$};
\node[blk] (rti) at (5.45,-1.0) {LPF};
\node[blk] (rtq) at (5.45,-2.0) {LPF};
\node[blk] (daci) at (6.85,-1.0) {DAC};
\node[blk] (dacq) at (6.85,-2.0) {DAC};
\node[dsp] (dtx) at (8.75,-1.5) {DUC, interp.\\CFR, DPD\\I/Q and LO\\calibration};
\draw[arr] (dtx.west |- daci.east) -- (daci.east); \draw[arr] (dtx.west |- dacq.east) -- (dacq.east);
\draw[arr] (daci)--(rti); \draw[arr] (dacq)--(rtq); \draw[arr] (rti)--(ti); \draw[arr] (rtq)--(tq);
\draw[arr] (ti.west) -| (sum.north); \draw[arr] (tq.west) -| (sum.south);
\draw[arr] (sum)--(pa); \draw[arr] (pa.west) -- (dup.east |- pa.west);
% frequency generation
\node[syn] (ref) at (2.65,0) {TCXO};
\node[syn] (syn) at (5.75,0) {PLL synthesisers, $\div2$ for 0$^\circ$/90$^\circ$};
\draw[lo] (ref) -- (syn);
\draw[lo] (4.1,0.31) -- node[right,font=\tiny,text=practc]{LO$_{\text{RX}}$} (mq.south);
\draw[lo] (4.1,-0.31) -- node[right,font=\tiny,text=practc]{LO$_{\text{TX}}$} (ti.north);
\draw[lo] (4.1,2.6) node[above,font=\tiny,text=practc]{LO 0$^\circ$} -- (mi.north);
\draw[lo] (4.1,-2.6) node[below,font=\tiny,text=practc]{LO 90$^\circ$} -- (tq.south);
\draw[arr,dashed] (dtx.north) -- node[right,font=\tiny,align=left]{observation\\receiver} (drx.south);
\end{tikzpicture}
\caption{A direct-conversion (zero-IF) transceiver, the architecture of nearly every integrated radio
chip including the AD9364 in the USRP B200. The receive path (top) amplifies, mixes to baseband in
quadrature, filters and digitises; the transmit path (bottom) mirrors it. Everything right of the
converters is DSP, including the correction of the analog impairments described in this chapter. The
dashed feedback path, an \emph{observation receiver}, lets the transmitter learn its own distortion.}
\label{fig:ch07_transceiver}
\end{figure}

\subsection{A whisper in a stadium}

Why is this hard? Look at Figure~\ref{fig:ch07_spectrum_mic}, a sketch of what a wideband antenna
in a city actually hears. The strongest signals---a nearby base station, a Wi-Fi router across the
desk---arrive at $-25$ to $-40$\,dBm. The thermal noise in a 200\,kHz channel is about $-121$\,dBm.
The signal you want might be at $-100$\,dBm, sitting 75\,dB below the loudest neighbour: a power ratio
of about thirty million. And GPS, the most-used radio service on Earth, arrives \emph{below} the noise
altogether (Chapter~\ref{ch:spreadspectrum}).

\mdcfig{ch07_spectrum_mic}{What a wideband antenna hears in a city (an illustrative sketch). The
receiver must detect a faint signal while strong neighbours, some only a few megahertz away, are
tens of decibels louder. Every block of the receiver is designed around this picture: the filters
remove what they can, the LNA lifts the whisper above its own hiss, and the linearity of everything
after it decides whether the shouts distort the whisper.}

\begin{analogy}[A microphone for the electromagnetic spectrum]
A receiver is a microphone pointed at the radio spectrum. A good studio microphone must do two
contradictory things: pick up a whisper without adding hiss (\emph{noise figure}), and survive a
drum kit without distorting (\emph{linearity}). It also must not hum (\emph{spurious signals}) and
must stay at the right pitch (\emph{frequency accuracy}). Radio engineers fight exactly the same four
battles. Where the analogy breaks: a microphone hears everything at once and your brain picks out the
whisper, whereas a radio must \emph{throw away} the shouts with filters before they reach the parts
that would distort, because no amount of clever processing afterwards can undo the damage.
\end{analogy}

Each block in Figure~\ref{fig:ch07_transceiver} contributes noise, distortion, spurious signals and
frequency errors, and the radio designer's job is to allocate a budget for each. The budget comes from
the standard the radio must meet: a reference sensitivity (the weakest signal it must decode), a set
of blocking and intermodulation tests (the strongest interferers it must tolerate while doing so), an
error-vector magnitude (EVM) for the transmitted signal, an adjacent-channel leakage ratio and a
spectrum emission mask (how much power it may spray onto its neighbours), and frequency-error limits.
Almost every number in this chapter is one of those requirements traced back to a circuit
specification.

\mdcphoto[0.85]{ch07_iphone_board}{The main board of an iPhone~5 (2012). The left third is radio:
power-amplifier and front-end modules, the transceiver chip and the filters that select dozens of
bands. The big chips in the middle and right are the processor and memory. A phone's radio is not one
chip but a small city of them, and Section~\ref{sec:ch07:frontend} explains why.}

\begin{keyidea}[Four quantities run the radio]
Strip a radio data sheet to its essentials and four numbers remain: the \emph{noise figure}, which sets
how weak a signal it can hear; the \emph{intercept points and compression point}, which set how strong
a signal it can tolerate; the \emph{phase noise} of its oscillators, which limits both the
constellation density it can support and its tolerance of nearby blockers; and its \emph{efficiency},
which sets battery life, heat and the electricity bill of a network. Every architecture in this chapter
is a different trade among those four, and every digital correction algorithm exists because one of
them could not be met in analog alone.
\end{keyidea}

%======================================================================
\section{Receiver architectures}
\label{sec:ch07:arch}
%======================================================================

\subsection{The mixer and the mirror-image station}

Every receiver architecture is an answer to one question: \emph{how do we move the signal from where
it is (say 2.4\,GHz) to where we can process it (near zero hertz)?} The tool is the mixer, a
\emph{frequency translator}. Multiply a signal at $f_{\text{RF}}$ by a local oscillator at
$f_{\text{LO}}$ and, by the product-to-sum identity, you get copies at $f_{\text{RF}}-f_{\text{LO}}$
and $f_{\text{RF}}+f_{\text{LO}}$ (Figure~\ref{fig:ch07_mixer_translator}). Keep the difference, filter
away the sum, and the signal---every bit of its modulation intact---now lives at a lower,
friendlier frequency.

\mdcfig{ch07_mixer_translator}{A mixer is a frequency translator. A 900\,MHz signal multiplied by an
830\,MHz local oscillator appears at the difference, 70\,MHz, and at the sum, 1730\,MHz. The modulation
moves unchanged; a filter keeps whichever copy we want.}

There is a catch, and it has haunted radio since Armstrong. The mixer cannot tell the difference
between a signal \emph{above} the LO and one \emph{below} it: both $f_{\text{LO}}+f_{\text{IF}}$ and
$f_{\text{LO}}-f_{\text{IF}}$ come out at the same $f_{\text{IF}}$. The second one is the \term{image}
(Figure~\ref{fig:ch07_image_mirror}), and once it has passed through the mixer it can never be
separated from the wanted signal again.

\begin{analogy}[A mirror-image station]
Stand the LO up as a mirror on the frequency axis. Every station on one side has a twin reflected on
the other, and the mixer's output is what you see in the mirror: the wanted station and its reflection
superimposed. The superheterodyne's preselector filter is a curtain drawn over one half of the room;
the image-reject mixer is a trick of polarised light that cancels the reflection; the zero-IF receiver
puts the mirror right through the middle of the wanted station, so the station becomes its own
reflection. Every architecture below is one of these three ideas.
\end{analogy}

\mdcfig{ch07_image_mirror}{The image frequency. A station mirrored about the LO, at
$f_{\text{LO}}-f_{\text{IF}}$, lands on exactly the same IF as the wanted station at
$f_{\text{LO}}+f_{\text{IF}}$. The two are $2f_{\text{IF}}$ apart at RF, and only a filter (or a
phasing trick) ahead of the mixer can tell them apart.}

\subsection{Four families}

Figure~\ref{fig:ch07_architectures} sketches the four architectures in use today in the frequency
domain. The \term{superheterodyne} of Chapter~\ref{ch:analog} converts to a fixed intermediate frequency
(IF) where high-quality filters select the channel. Its weakness is the image: a mixer with LO
at $f_{\text{LO}}$ converts both $f_{\text{LO}}+f_{\text{IF}}$ and $f_{\text{LO}}-f_{\text{IF}}$ to the
same IF, so a filter \emph{before} the mixer must remove whichever of the two is unwanted. A high IF
makes the image easy to filter (it is $2f_{\text{IF}}$ away) but the channel hard to select; a low IF
does the reverse; double and triple conversion resolve the conflict at the cost of more mixers,
oscillators and spurious responses. The superhet remains the choice where selectivity and dynamic range
are paramount: spectrum analysers, military and surveillance receivers, many base stations, broadcast
and satellite ground equipment, and the RF sections of high-end amateur transceivers.

The \term{direct-conversion} or \term{zero-IF} receiver sets the LO at the carrier and produces I and
Q baseband outputs directly. It needs no image filter and no IF filter---channel selection is a
low-pass filter that can be integrated on chip---so a single die can cover 70\,MHz to 6\,GHz, as the
AD9364 in the B200 does. That flexibility is why zero-IF dominates mobile phones, integrated
transceivers and SDRs. How does it escape the image? With two mixers driven by LOs $90^\circ$ apart,
it keeps both the I and the Q of the signal, and so knows which side of the LO each component came
from: positive and negative frequencies are distinguishable in a complex signal
(Chapter~\ref{ch:signals}). Its weaknesses all come from putting the wanted signal at DC, and the next
subsection examines them one by one.

\mdcfig{ch07_architectures}{Receiver architectures in the frequency domain. Top: the superheterodyne
must reject the image with an RF filter before mixing. Second: the zero-IF receiver's ``image'' is the
signal's own mirror, rejected only by I/Q balance, and its baseband sits on the DC offset. Third: a
low-IF receiver keeps away from DC but its image is a neighbouring channel. Bottom: a direct
RF-sampling receiver digitises the whole band and performs the mixing digitally.}

A \term{low-IF} receiver places the LO about half a channel (or a little more) away from the carrier, so
the wanted signal lands at a small IF, typically a few hundred kilohertz to a few megahertz, avoiding
the DC and flicker problems while keeping channel filtering on chip. The image is now the
\emph{adjacent} channel, which may be much stronger than the wanted one, and must be rejected by I/Q
balance (Lab~15 measures exactly this trade). Low-IF is common in Bluetooth, GNSS, FM-radio and
narrowband IoT chips.

\term{Direct RF sampling} puts the ADC at the antenna, or nearly so. Converters sampling at several
gigasamples per second with 12 to 14 bits, integrated with FPGA logic in devices such as AMD's
Zynq UltraScale+ RFSoC, can digitise everything up to several gigahertz; the mixers, IF filters and
channel filters all become DSP (Chapter~\ref{ch:dsp}). This is the logical end of Mitola's software radio
vision, limited today by power consumption, by the clock jitter that limits the SNR of a converter
sampling a high-frequency signal (Chapter~\ref{ch:pcm}), and by the dynamic range needed when strong and
weak signals share the band.

\begin{tryit}[Hear the image, then kill it]
Open Lab 15 (\lab{lab15\_superhet\_receiver.py}) and choose \emph{The image problem}. Set the wanted
station to 1000\,kHz with a 455\,kHz IF, then drag \emph{Image station strength} up to 40\,dB and
listen to the mirror station break through. Now raise \emph{Loaded Q} and the number of
\emph{Tuned circuits}, or switch the \emph{Intermediate frequency} to a higher value, and watch the
image fall away. Finally open \emph{Zero-IF vs low-IF} and compare the two architectures with the same
\emph{Gain imbalance} and \emph{Phase error}.
\end{tryit}

\subsection{The trouble with zero: direct-conversion impairments}
\label{sec:ch07:zeroif}

Putting the wanted signal at 0\,Hz is like parking your car exactly where everyone else dumps their
rubbish. Figure~\ref{fig:ch07_zeroif_impairments} is a simulation of what a direct-conversion receiver
hands its ADCs, with every impairment exaggerated just enough to be visible. There are five.

\mdcfig{ch07_zeroif_impairments}{Simulated spectrum at the output of a zero-IF receiver's analog
section, before digital channel filtering. A 16-QAM signal sits at DC; an amplitude-varying blocker
40\,dB stronger sits at $+6$\,MHz. The receiver adds a DC offset from LO self-mixing, flicker noise
and the blocker's second-order intermodulation (its envelope) near DC, and an image of the blocker
at $-6$\,MHz from I/Q imbalance.}

\paragraph{DC offset and LO self-mixing.} The LO signal is many decibels stronger than anything at
the antenna, and some of it leaks: through substrate and package coupling to the LNA input, or out of
the antenna and back off a nearby reflector. Leaked LO arriving at the mixer's RF port mixes with the
LO itself, $\cos^2(\omega_{\text{LO}}t)=\tfrac12+\tfrac12\cos2\omega_{\text{LO}}t$, producing a DC
term. A leakage of $-60$\,dBm referred to the LNA input, multiplied by 60\,dB of receiver gain, is a
0\,dBm DC offset at the ADC---larger than the wanted signal and possibly large enough to saturate the
baseband amplifiers. Worse, the offset changes with LO frequency, gain setting, temperature and the
environment around the antenna (a hand moving near a phone changes the reflection), so a one-time
calibration is not enough. Remedies: a careful layout and LO-buffer design that minimise the leakage;
DC-offset cancellation loops in the analog baseband that run at power-up and at every gain change;
a digital high-pass filter or running-mean subtraction after the ADC (Section~\ref{sec:ch07:iq}); and,
for signals whose spectrum has a hole at DC (OFDM with an unused DC subcarrier), simply ignoring it.

\paragraph{LO radiation.} The same leakage path running outward means a zero-IF receiver radiates its
LO \emph{in the band it is receiving}, where other users can be disturbed. Standards therefore limit
spurious emissions in receive mode, and the LO is often generated at twice or two-thirds the carrier
and divided down, so that the VCO itself is not at the antenna frequency.

\paragraph{Flicker noise.} Baseband amplifiers in CMOS have $1/f$ noise corners from tens of
kilohertz to megahertz (Chapter~\ref{ch:noise}). In a superhet that noise lies far from the IF; in a
zero-IF receiver it falls right on the centre of the wanted channel. For wideband signals---a 20\,MHz
LTE carrier, say---the few hundred kilohertz around DC carry a small fraction of the energy and the
penalty is modest; for a 200\,kHz GSM channel it is serious, which is why early GSM receivers stayed
superheterodyne and why narrowband receivers prefer low-IF.

\paragraph{Second-order intermodulation.} A mixer is a deliberately nonlinear device, and its
even-order distortion, $a_2x^2$, demodulates the \emph{envelope} of any strong signal. A blocker
$b(t)\cos(\omega_bt+\theta(t))$ produces $\tfrac12a_2b^2(t)$, a baseband term with roughly twice the
blocker's envelope bandwidth, centred at DC, whatever the blocker's frequency. A constant-envelope
blocker produces only a DC shift; an AM, OFDM or pulsed blocker (a GSM burst, a TDD neighbour switching
on and off) produces a noise-like interferer right on top of the wanted signal. The figure of merit is
the \term{second-order intercept point} (IIP2), and zero-IF receivers need IIP2 of $+45$ to $+60$\,dBm,
attained with balanced, carefully matched mixers and often with on-chip IIP2 calibration. The most
dangerous case is a full-duplex FDD phone: its own transmitter, leaking through the duplexer at perhaps
$-25$\,dBm, is an amplitude-modulated blocker whose IM2 lands on the receive channel.
Figure~\ref{fig:ch07_im2_envelope} shows the mechanism: square a signal that lives far away in
frequency, and its slowly varying envelope appears at baseband, as if the mixer had become an AM
detector.

\mdcfig{ch07_im2_envelope}{How IM2 brings a distant blocker to DC. Left: a strong blocker with a
varying envelope, many megahertz from the wanted channel. Right: after a square-law term and
low-pass filtering, what remains is the blocker's envelope---a noise-like interferer sitting right on
top of a zero-IF receiver's wanted signal, whatever the blocker's frequency.}

\paragraph{I/Q imbalance.} If the I and Q paths differ slightly in gain or the two LO phases are not
exactly $90^\circ$ apart, every signal at baseband frequency $f$ acquires an image at $-f$. In a zero-IF
receiver the image of the wanted signal lands on itself (its mirror half), which merely distorts the
constellation; but a strong signal at $+f$ elsewhere in the captured band sends its image into whatever
is at $-f$. Section~\ref{sec:ch07:iq} derives the image-rejection ratio and the algorithms that correct
it.

\begin{tryit}[Fingerprints of a zero-IF radio]
In Lab 1 (\lab{lab01\_baseband.py}), open \emph{Zero-IF fingerprints: DC and I/Q images}. Raise the
\emph{DC offset} slider and watch the spike at the centre of the spectrum; then add \emph{Gain
imbalance} and \emph{Phase imbalance} and move \emph{Your channel offset} until a strong
\emph{Mirror neighbour} lands on top of your channel. Toggle \emph{Remove DC} and \emph{Blind I/Q
correction} to see the two cures of Section~\ref{sec:ch07:iq} at work.
\end{tryit}

\begin{history}[The zero-IF receiver comes in from the cold]
Direct conversion is as old as radio---the homodyne and the synchrodyne were described between the
1920s and 1940s, and radio amateurs have built direct-conversion receivers for CW and SSB for decades---but
its integration began with pagers. In the early 1980s Ian Vance and colleagues at STC in the UK
built a pager receiver whose channel selection was done entirely by on-chip low-pass filters, with the
signal at zero IF; the FSK paging signal had little energy at DC, so the DC offset could simply be
removed. Through the 1990s, as CMOS improved and DC-offset cancellation loops matured, the architecture
spread to DECT, to GSM and then to essentially all cellular, Wi-Fi and Bluetooth chips. Asad Abidi's
1995 paper ``Direct-conversion radio transceivers for digital communications'' laid out the problems
listed in this section so clearly that it remains the place to start. The superheterodyne did not die;
it retreated to where its selectivity and dynamic range are worth the filters.
\end{history}

\subsection{Image-reject mixers: Hartley and Weaver}
\label{sec:ch07:imagereject}

\begin{figure}[htbp]\centering
\begin{tikzpicture}[node distance=0.45cm and 0.5cm,
  blk/.style={draw=navy,fill=navy!6,thick,minimum height=0.7cm,minimum width=1.1cm,align=center,font=\sffamily\scriptsize},
  mix/.style={draw=navy,thick,circle,minimum size=0.55cm,inner sep=0pt,font=\sffamily\small},
  arr/.style={-{Stealth[length=1.6mm]},thick,navy}]
% Hartley
\node[font=\sffamily\scriptsize] (in) {RF};
\node[mix,right=0.6cm of in,yshift=0.6cm] (m1) {$\times$};
\node[mix,right=0.6cm of in,yshift=-0.6cm] (m2) {$\times$};
\node[blk,right=of m1] (l1) {LPF};
\node[blk,right=of m2] (l2) {LPF};
\node[blk,right=of l2,fill=accent!6,draw=accent] (h) {$90^\circ$\\shift};
\node[mix,right=2.5cm of l1,yshift=-0.6cm] (s) {$+$};
\node[right=0.4cm of s,font=\sffamily\scriptsize] (out) {IF, image cancelled};
\node[above=0.12cm of m1,font=\tiny] {$\cos\omega_{\text{LO}}t$};
\node[below=0.12cm of m2,font=\tiny] {$\sin\omega_{\text{LO}}t$};
\draw[arr] (in) -- ++(0.3,0) |- (m1); \draw[arr] (in) -- ++(0.3,0) |- (m2);
\draw[arr] (m1)--(l1); \draw[arr] (m2)--(l2); \draw[arr] (l2)--(h);
\draw[arr] (l1) -| (s); \draw[arr] (h) -| (s); \draw[arr] (s)--(out);
\node[left=0.1cm of in,font=\sffamily\scriptsize\bfseries,text=navy] {Hartley};
% Weaver
\node[font=\sffamily\scriptsize,below=2.7cm of in] (in2) {RF};
\node[mix,right=0.6cm of in2,yshift=0.6cm] (w1) {$\times$};
\node[mix,right=0.6cm of in2,yshift=-0.6cm] (w2) {$\times$};
\node[blk,right=of w1] (wl1) {LPF};
\node[blk,right=of w2] (wl2) {LPF};
\node[mix,right=of wl1] (w3) {$\times$};
\node[mix,right=of wl2] (w4) {$\times$};
\node[mix,right=0.7cm of w3,yshift=-0.6cm] (ws) {$-$};
\node[right=0.4cm of ws,font=\sffamily\scriptsize] (out2) {IF, image cancelled};
\node[above=0.12cm of w1,font=\tiny] {$\cos\omega_1t$};
\node[below=0.12cm of w2,font=\tiny] {$\sin\omega_1t$};
\node[above=0.12cm of w3,font=\tiny] {$\cos\omega_2t$};
\node[below=0.12cm of w4,font=\tiny] {$\sin\omega_2t$};
\draw[arr] (in2) -- ++(0.3,0) |- (w1); \draw[arr] (in2) -- ++(0.3,0) |- (w2);
\draw[arr] (w1)--(wl1); \draw[arr] (w2)--(wl2); \draw[arr] (wl1)--(w3); \draw[arr] (wl2)--(w4);
\draw[arr] (w3) -| (ws); \draw[arr] (w4) -| (ws); \draw[arr] (ws)--(out2);
\node[left=0.1cm of in2,font=\sffamily\scriptsize\bfseries,text=navy] {Weaver};
\end{tikzpicture}
\caption{Image-reject receivers. The Hartley architecture (top) cancels the image by a $90^\circ$
shift of one quadrature branch before summing; the Weaver architecture (bottom) achieves the same with a
second quadrature mixing stage. Either cancels the image only as well as its two branches match.}
\label{fig:ch07_hartley}
\end{figure}

There is a third way to deal with the image: cancel it by phasing, exactly as the phasing SSB modulator
of Chapter~\ref{ch:analog} cancels the unwanted sideband. It is the ``polarised light'' of our mirror
analogy. The \term{Hartley architecture} (Figure~\ref{fig:ch07_hartley}, top) mixes the RF input with
quadrature LOs, low-pass filters both outputs, shifts one of them by a further $90^\circ$ and adds. To
see why it works, write the input as a wanted component at $\omega_{\text{LO}}+\omega_{\text{IF}}$ and
an image at $\omega_{\text{LO}}-\omega_{\text{IF}}$, both of amplitude 1. After mixing with
$\cos\omega_{\text{LO}}t$ and $\sin\omega_{\text{LO}}t$ and low-pass filtering,
\begin{align}
I(t)&=\tfrac12\cos\omega_{\text{IF}}t+\tfrac12\cos\omega_{\text{IF}}t, \nonumber\\
Q(t)&=-\tfrac12\sin\omega_{\text{IF}}t+\tfrac12\sin\omega_{\text{IF}}t,
\end{align}
where in each line the first term is the wanted signal and the second the image. Shifting $Q$ by
$90^\circ$ turns $\mp\sin$ into $\pm\cos$: the wanted terms add and the image terms cancel. In complex
terms, the I/Q pair is the complex envelope, in which the wanted signal sits at $+\omega_{\text{IF}}$ and
the image at $-\omega_{\text{IF}}$, and the Hilbert transform plus addition is simply ``take the
positive-frequency half''. A low-IF receiver with digital I/Q processing is a Hartley receiver whose
$90^\circ$ shift and sum happen in DSP. The image rejection is limited by the same gain and phase
matching as any I/Q system: about 30 to 40\,dB from good analog design, more with calibration.

The \term{Weaver architecture} (bottom of the figure) replaces the wideband $90^\circ$ phase shifter,
which is hard to build accurately over a range of IF frequencies, with a second quadrature mixing stage.
It is the Weaver SSB method of Chapter~\ref{ch:analog} turned into a receiver, and it is the structure of
every receiver that performs a quadrature analog mix followed by a complex digital mix.

\subsection{Direct RF sampling and the RFSoC}

The direct-sampling receiver is easiest to understand as a bandpass sampler (Chapter~\ref{ch:pcm}):
an ADC sampling at $f_s$ folds every Nyquist zone of width $f_s/2$ onto the first, so a band-pass filter
selecting one zone, followed by a fast ADC and a digital down-converter, is a complete receiver. Modern
RF data converters integrated with FPGA fabric---the Gen~3 RFSoC devices sample at several GS/s with
14-bit resolution and accept inputs up to roughly 6\,GHz, according to AMD's data sheets---make this
practical for base stations, phased-array radars, test equipment and the high-end SDRs of
Section~\ref{sec:ch07:platforms}. The trade-offs are power (a multi-GS/s converter and the DSP behind it
burn watts, not milliwatts), the converter's own noise figure (Chapter~\ref{ch:noise} shows that a
converter's noise density makes it a poor front-end unless preceded by an LNA with gain), clock jitter,
and the spurious products of the converter's interleaved cores, which appear at predictable offsets and
must be calibrated out.

Table~\ref{tab:ch07:arch} summarises the four. If you remember only one row, remember the second:
it is the architecture inside your phone and inside your B200.

\begin{table}[htbp]\centering\small
\caption{Receiver architectures compared.}
\label{tab:ch07:arch}
\begin{tabular}{>{\raggedright\arraybackslash}p{2.3cm}>{\raggedright\arraybackslash}p{3.3cm}>{\raggedright\arraybackslash}p{3.4cm}>{\raggedright\arraybackslash}p{3.4cm}}
\toprule
Architecture & Strengths & Weaknesses & Typical use \\
\midrule
Superhet\-ero\-dyne & Best selectivity and dynamic range; no DC problems & Image and IF filters (off chip), spurious responses, cost & Spectrum analysers, military, satellite and broadcast receivers \\
Zero-IF & Fully integrable; one design covers many bands; simple & DC offset, LO leakage, flicker noise, IM2, I/Q imbalance & Phones, Wi-Fi, SDR RFICs (AD9364, AD9361, LMS7002M) \\
Low-IF & Integrable; away from DC and $1/f$ noise & Image is adjacent channel; needs high IRR & Bluetooth, GNSS, FM, IoT \\
Direct RF sampling & No analog mixing; whole band at once & Power, ADC noise and spurs, jitter & Base stations, radar, test gear, RFSoC SDRs \\
\bottomrule
\end{tabular}
\end{table}

%======================================================================
\section{The RF front end}
\label{sec:ch07:frontend}
%======================================================================

Open any phone and the radio parts that look most like ``old-fashioned electronics''---little metal
cans, ceramic blocks, tiny packages with no logic in them at all---are the front end. Everything
between the antenna connector and the mixer is where most of a receiver's noise figure and much of its
linearity are decided, and where the hardest physical limits live. It consists of filters, switches, a
duplexer if the radio is full-duplex, and the LNA.

\mdcphoto[0.72]{ch07_saw}{Inside a surface-acoustic-wave filter (a Siemens IF filter for television
receivers). The comb-like interdigital transducers on the piezoelectric crystal launch and catch an
acoustic wave; the spacing and length of the fingers set the frequency response. Modern RF SAW filters
in phones work the same way at gigahertz frequencies, with fingers well under a micrometre wide.}

\subsection{Filters: SAW, BAW, ceramic and cavity}

A front-end filter must have low insertion loss in the pass band (every decibel adds directly to the
noise figure, Chapter~\ref{ch:noise}) and steep skirts to reject strong signals just outside the band,
all in a package smaller than a grain of rice. The trick that makes this possible is wonderfully
indirect: \emph{turn the radio wave into sound}. An acoustic wave travels about a hundred thousand
times slower than an electromagnetic one, so its wavelength at 2\,GHz is a couple of micrometres
instead of fifteen centimetres, and a resonator that would be a hand-sized cavity becomes a pattern
on a chip (Figure~\ref{fig:ch07_saw}). Four technologies cover the range.
\begin{description}
\item[Surface acoustic wave (SAW)] filters convert the electrical signal into an acoustic wave on a
piezoelectric substrate (lithium tantalate or niobate) with interdigital transducers, whose geometry
sets the response. They are tiny and cheap and dominate below about 2.5\,GHz; temperature-compensated
(TC-SAW) variants hold their edges over temperature for tightly spaced duplex bands.
\item[Bulk acoustic wave (BAW)] filters, including the film bulk acoustic resonator (FBAR), use a
thin piezoelectric film (aluminium nitride, often scandium-doped) resonating through its thickness.
They have higher $Q$ at higher frequencies and handle more power, and they dominate the 2 to 7\,GHz
bands of LTE, 5G NR and Wi-Fi~6E/7.
\item[Ceramic and LTCC] filters and resonators are used where size matters less: GNSS, ISM-band
modules, and the band-select filters of simple radios.
\item[Cavity and dielectric-resonator] filters, with unloaded $Q$ in the thousands, are used in base
stations, where they must pass hundreds of watts and reject the receive band by 80 to 100\,dB, and in
satellite payloads and ground stations.
\end{description}

\subsection{Duplexers and the FDD problem}

A frequency-division duplex radio (Chapter~\ref{ch:multipleaccess}) transmits and receives at the same
time on two frequencies separated by the duplex spacing. It is like trying to listen for a whisper
while you yourself are shouting---the shout is only a little way off in pitch. A \term{duplexer} is two
band-pass filters joined at the antenna: the transmit filter passes the uplink band and blocks the
downlink, the receive filter does the reverse (Figure~\ref{fig:ch07_duplexer}).

\mdcfig{ch07_duplexer}{An LTE band-3 duplexer (illustrative filter shapes). The transmit filter passes
the 1710--1785\,MHz uplink, the receive filter the 1805--1880\,MHz downlink, only 20\,MHz above. Even with
55\,dB of isolation the phone's own $+23$\,dBm transmitter arrives at its LNA at $-32$\,dBm---far
stronger than anything else the receiver will see.}

Consider an LTE phone transmitting $+23$\,dBm in band~3 (uplink 1710--1785\,MHz,
downlink 1805--1880\,MHz) while trying to receive a signal near $-100$\,dBm only 95\,MHz away. If the
duplexer's transmit-to-receive isolation is 55\,dB, the transmitter's carrier arrives at the LNA at
$-32$\,dBm---as a blocker that compresses the LNA, generates IM2 in a zero-IF mixer, and mixes with
phase noise. Just as dangerous is the transmitter's own \emph{noise} in the receive band: a PA whose
output noise density at the receive frequency is $-130$\,dBm/Hz would, after 55\,dB of isolation, add
$-185$\,dBm/Hz, harmlessly below the thermal floor; but a cheaper transmit chain with $-120$\,dBm/Hz would
add $-175$\,dBm/Hz, degrading sensitivity by about 2.5\,dB. That is why the duplexer specification has
two isolation numbers (in the transmit band and in the receive band), and why the transmit chain's noise
at the receive frequency is a specification of its own.

\begin{inpractice}[The band-count problem in a smartphone]
A 2G phone of 1995 needed two or three bands. A flagship 5G phone today supports several dozen LTE and
NR bands from 600\,MHz to 6\,GHz (plus millimetre-wave modules), with carrier aggregation that requires
several receive and transmit chains to operate simultaneously on different bands. Each FDD band needs
its own duplexer; carrier aggregation needs multiplexers (quadplexers, hexaplexers) whose filters do not
load one another; and every combination of simultaneously active bands creates new intermodulation and
harmonic hazards (a band-12 uplink's third harmonic, for instance, must not land in a band being
received). The result is that the RF front end of a modern phone contains of the order of a hundred
acoustic filters, packed with switches, LNAs and PAs into a handful of front-end modules made by a small
number of specialised companies. Filter count, not transistor count, is the scaling problem of the
mobile radio, and it is one reason why tunable filters and full-duplex cancellation
(Section~\ref{sec:ch07:fullduplex}) are active research areas.
\end{inpractice}

\subsection{The low-noise amplifier}

Remember the party game in which a whispered message passes down a line of people? The person who
matters most is the \emph{first} one. If they mishear, nobody after them can recover the message; if
they hear it clearly and repeat it loudly, the rest of the line hardly matters. The LNA is that first
person in the whisper chain, and Friis's formula (equation~\eqref{eq:ch03:friis}) is the mathematics of
the game.

The LNA has three jobs: provide enough gain to make the noise of everything after it negligible; add
as little noise as possible itself; and remain linear in the presence of the strongest signals the band
filter lets through. These are in tension. Figure~\ref{fig:ch07_whisper_chain} shows who adds the
noise in the line-up we shall analyse in Section~\ref{sec:ch07:metrics}: the lossy switch in front and
the LNA itself account for nearly all of it, and the mixer and baseband, for all their 10 and 20\,dB
noise figures, barely register.

\mdcfig[0.5]{ch07_whisper_chain}{Who adds the noise? Contributions $(F_k-1)/G_{\text{before}}$ of
each stage of the worked line-up of Section~\ref{sec:ch07:metrics}. The passive loss before the LNA
costs as much as the LNA itself; the noisy mixer and baseband barely matter behind 17.5\,dB of gain.}

\begin{analogy}[The first person in a whisper chain]
In a whisper chain, the first listener's mistakes are copied by everyone; later listeners hear a loud,
clear version and add little. In a receiver, the first stage's noise is amplified by every later
stage, while a later stage's noise is divided by all the gain in front of it. So spend your best
ears---the lowest noise figure---at the front, and make the first person speak up (gain) so the rest of
the line cannot spoil the message. The analogy's limit: shout too loudly and the next listener's ears
ring (they compress), which is why the LNA's gain cannot simply be made huge.
\end{analogy}

\paragraph{Noise match versus power match.} Chapter~\ref{ch:noise} showed that a transistor's noise
factor depends on its source admittance, $F=F_{\min}+(R_n/G_s)|Y_s-Y_{\text{opt}}|^2$
(equation~\eqref{eq:ch03:fmin}), and that the optimum source for noise, $Y_{\text{opt}}$, is not in
general the conjugate match that maximises gain. Presenting $Y_{\text{opt}}$ gives the lowest noise
figure but a poor input return loss: power is reflected back toward the filter, whose response is
designed for a 50\,$\Omega$ load and becomes distorted when mismatched. The classic resolution in CMOS
is \emph{inductive source degeneration}: an inductor $L_s$ in the source creates a real input resistance
$g_mL_s/C_{gs}=\omega_TL_s$ without adding thermal noise, and with a series gate inductor resonating the
input capacitance, the power match and noise match can be brought close together at one frequency.
Wideband LNAs (for SDR front ends covering decades of frequency) instead use resistive feedback or
noise-cancelling topologies, accepting a somewhat higher noise figure for a flat match.

\paragraph{Gain, noise figure and linearity.} A typical integrated cellular LNA has a gain of 15 to
20\,dB, a noise figure of 1 to 2\,dB and an IIP3 of $-10$ to 0\,dBm; a discrete GaAs pHEMT LNA for a
satellite ground terminal may reach a noise figure of 0.3\,dB at 4\,GHz. High gain helps the cascade
noise figure but pushes the mixer toward overload. Real receivers therefore provide \emph{gain steps}:
the LNA can be switched to a lower gain or bypassed altogether when the received signal is strong,
which is the first and coarsest stage of the automatic gain control of Section~\ref{sec:ch07:agc}.
The AD9364's receive gain table, which UHD exposes as a single number from 0 to about 76\,dB on the B200,
is a sequence of such settings across LNA, mixer, TIA and baseband amplifiers, chosen to keep the noise
figure low at high gain and the linearity high at low gain.

\mdcfig[0.5]{ch07_cable_loss}{Why the LNA belongs at the antenna. Every decibel of cable loss in
front of an SDR (red) adds a decibel of noise figure; put a 20\,dB, 0.8\,dB-NF preamplifier at the
antenna (green) and the same cable costs almost nothing. Cable loss is approximate, for thin coax at
2.4\,GHz.}

\begin{pitfall}[Loss before the LNA counts double]
A decibel of loss ahead of the LNA---a long coaxial cable, a switch, a lossy filter, a bad connector---adds
a full decibel to the system noise figure, because it attenuates the signal while the noise at its output
is still the thermal $kT$ (Chapter~\ref{ch:noise}). On the bench, a B200 connected to an antenna through
5\,m of thin RG-174 at 2.4\,GHz has lost several decibels of sensitivity before the signal reaches the
board. That is why satellite and radio-astronomy receivers mount the LNA at the antenna feed, why tower-mounted
amplifiers exist for base stations, and why an external low-noise preamplifier at the antenna is the
cheapest sensitivity upgrade for any SDR.
\end{pitfall}

\mdcfig{ch07_switching_mixer}{A switching mixer at work. The RF signal (top) is multiplied by a
$\pm1$ square wave at the LO frequency (middle). The product (bottom, grey) looks like a mess, but a
low-pass filter pulls out the difference frequency (red)---the translated signal, here at one unit.}

\subsection{Mixers}

A mixer multiplies the RF signal by the LO. Since a perfect multiplier is hard to build, practical
mixers \emph{switch}: they multiply the RF signal by a square wave at the LO frequency, flipping the
signal's polarity on and off at the LO rate. A square wave contains the fundamental plus odd harmonics
at $3f_{\text{LO}}$, $5f_{\text{LO}}$ and so on, and that, as we shall see, has consequences
(Figure~\ref{fig:ch07_switching_mixer}).

\paragraph{Passive mixers.} The diode ring of Chapter~\ref{ch:analog} and its CMOS successor, a ring of
four MOS switches driven by a large LO, are \emph{passive}: they have conversion loss rather than gain
(about 6\,dB for an ideal switching mixer with a resistive load), but excellent linearity (IIP3 of
$+10$ to $+30$\,dBm) and very little flicker noise, since no DC current flows through the switches. Modern
zero-IF receivers almost universally use a passive current-mode mixer: the LNA delivers a current, the
switches commutate it, and a transimpedance amplifier (TIA) at baseband converts it to a voltage while
presenting a low impedance that keeps voltage swings---and hence distortion---small at the mixer.

\paragraph{Active mixers and the Gilbert cell.} The \term{Gilbert cell} uses a differential pair, driven
by the RF signal, whose tail currents are steered by two cross-coupled differential pairs switched by
the LO. It provides conversion gain (10\,dB or so), good port-to-port isolation from its balanced
structure, and needs only a modest LO drive; its price is a higher noise figure (8 to 12\,dB SSB is
typical), flicker noise at the output (which matters at zero IF) and lower linearity than a passive
ring. It was the workhorse of bipolar radio ICs for thirty years.

\paragraph{Conversion gain, noise figure and LO drive.} A mixer is specified by its \emph{conversion gain}
(output power at the IF over input power at RF), its noise figure---single-sideband or double-sideband,
a factor of 2 apart, as Chapter~\ref{ch:noise} warned---its IIP3 and IIP2, its \emph{port isolations}
(LO-to-RF, which sets LO leakage and hence DC offset; LO-to-IF; RF-to-IF), and the LO power it needs. A
passive diode mixer needs $+7$ to $+17$\,dBm of LO (the ``level 7'' and ``level 17'' mixers of
catalogue fame); an integrated CMOS switching mixer needs rail-to-rail square-wave LO drive, generated
on chip by dividers and buffers that themselves consume a good fraction of the receiver's power.

\mdcfig[0.5]{ch07_harmonic_mixing}{Harmonic responses of a switching mixer. A square-wave LO has odd
harmonics at $1/n$ of the fundamental, so a receiver tuned to 100\,MHz with no preselector also
converts 300\,MHz with only 9.5\,dB less gain, 500\,MHz with 14\,dB less, and so on.}

\paragraph{Harmonic responses.} Because a switching mixer multiplies by a square wave, it responds at
$3f_{\text{LO}}$ with only about 10\,dB less conversion gain than at $f_{\text{LO}}$ (the square wave's
third harmonic is $1/3$ of its fundamental), at $5f_{\text{LO}}$ about 14\,dB less, and so on
(Figure~\ref{fig:ch07_harmonic_mixing}). A narrowband receiver with a band filter in front does not
care. A wideband SDR tuned to 100\,MHz with no preselector, however, will also receive signals at
300\,MHz and 500\,MHz with modest attenuation---a common source of ``mysterious'' signals in wideband
SDR captures. If you tune a B200 to an empty part of the FM band and see a station that should not be
there, ask yourself what lives at three times the frequency. \emph{Harmonic-rejection mixers}, which
combine several phases of the LO weighted to cancel the third and fifth harmonics, and switched
preselection filter banks (the B200 has none; higher-end USRPs do) are the remedies. Until then, a
simple external band-pass or low-pass filter in front of the SDR's input is the cheapest cure for
ghosts.

\Needspace{8\baselineskip}
\begin{history}[Barrie Gilbert's multiplier]
In 1968 Barrie Gilbert, then at Tektronix, published ``A precise four-quadrant multiplier with
subnanosecond response'' in the \emph{IEEE Journal of Solid-State Circuits}. His circuit---three
differential pairs and a translinear input stage---made analog multiplication accurate, wideband and
integrable. Gilbert noted later that a similar topology had been used earlier by Howard Jones for a
synchronous detector, but it was Gilbert's analysis of the translinear principle that made the circuit
a building block. The ``Gilbert-cell mixer'' became, through the 1980s and 1990s, the standard mixer
of every integrated receiver, from pagers to the first generations of GSM and Wi-Fi chips, and Gilbert
spent the rest of his career at Analog Devices, the company that would later build the RFICs at the heart
of today's SDRs.
\end{history}

%======================================================================
\section{Receiver figures of merit}
\label{sec:ch07:metrics}
%======================================================================

Every receiver data sheet is a list of compromises written as numbers. This section teaches you to
read them---and, better, to compute them for a receiver of your own design, stage by stage, the way
RF engineers do on the back of an envelope (or, these days, in a spreadsheet).

\subsection{Sensitivity and the noise-figure cascade}

The \term{sensitivity} of a receiver is the weakest signal it can demodulate at a specified error
rate. Chapter~\ref{ch:noise} derived it, equation~\eqref{eq:ch03:sens}; in words, start from the thermal
floor, add the bandwidth you let in, add the noise the receiver adds, and add the SNR the demodulator
needs:
\begin{equation}
P_{\text{sens}}=-174\,\text{dBm/Hz}+10\log_{10}B+\text{NF}+\text{SNR}_{\text{min}}\quad\text{dBm},
\label{eq:ch07:sens}
\end{equation}
and the Friis formula, equation~\eqref{eq:ch03:friis}, gives the noise figure of a cascade,
$F=F_1+(F_2-1)/G_1+(F_3-1)/(G_1G_2)+\cdots$. In a well-designed receiver the first one or two stages set
the noise figure and the later ones contribute a few tenths of a decibel.

\mdcfig{ch07_lineup_cumulative}{The worked line-up below, stage by stage. Left: the cumulative noise
figure jumps at the lossy switch and the LNA and then barely moves. Right: the cumulative IIP3 falls
with every stage of gain, and the last stage sets it. Noise is decided at the front, linearity at the
back.}

\begin{worked}[A B200-like receiver line-up]
Consider a receiver modelled loosely on an SDR with an integrated zero-IF RFIC. The numbers are
illustrative, chosen to be typical of such a chain, not taken from a data sheet:
\begin{center}\small
\begin{tabular}{lrrr|rrr}
\toprule
Stage & Gain & NF & IIP3 & Cum.\ gain & Cum.\ NF & Cum.\ IIP3 \\
 & (dB) & (dB) & (dBm) & (dB) & (dB) & (dBm) \\
\midrule
T/R switch, balun, traces & $-2.5$ & 2.5 & --- & $-2.5$ & 2.5 & --- \\
LNA & 20 & 1.5 & $-10$ & 17.5 & 4.0 & $-7.5$ \\
Mixer + TIA & 15 & 10 & $+5$ & 32.5 & 4.27 & $-13.7$ \\
Baseband filter + VGA & 30 & 20 & $+15$ & 62.5 & 4.36 & $-19.0$ \\
\bottomrule
\end{tabular}
\end{center}
\emph{Noise figure.} The 2.5\,dB of passive loss is a noise factor of 1.778 and a gain of 0.562. The LNA
adds $(1.413-1)/0.562=0.734$, giving $F=2.51$ (4.0\,dB): the loss ahead of the LNA has cost more than the
LNA itself. The mixer adds $(10-1)/(0.562\times100)=0.160$ and the baseband $(100-1)/1778=0.056$, for a
total of $F=2.73$, NF $=4.36$\,dB. The 20\,dB-noise-figure baseband stage contributes 0.09\,dB.

\emph{Linearity.} For IIP3 the cascade formula (derived below) adds the reciprocals of each stage's IIP3
referred to the input, $G_{\text{before}}/\text{IIP3}_k$, in milliwatts: $0.562/0.1=5.62$ for the LNA,
$56.2/3.16=17.8$ for the mixer and $1778/31.6=56.2$ for the baseband, a total of 79.6, so
$\text{IIP3}=1/79.6\,\text{mW}=-19.0$\,dBm. The \emph{last} stage, with the highest IIP3 but the most
gain in front of it, dominates.

\emph{Dynamic range.} In 1\,MHz the noise floor is $-174+60+4.4=-109.6$\,dBm and, by
equation~\eqref{eq:ch07:sfdr} below, the SFDR is $\tfrac23(-19.0+109.6)=60$\,dB. These values---NF of a
few decibels, IIP3 near $-20$\,dBm at maximum gain---are of the order Ettus quotes for the B200 (noise
figure below about 8\,dB including the board, IIP3 about $-20$\,dBm at typical noise figure). Reducing
the receive gain by 10\,dB in the baseband VGA would raise the cascade IIP3 to about $-15$\,dBm (the
VGA's output intercept stays put, so its input intercept rises by 10\,dB and the mixer now dominates) at
a cost of only a few tenths of a decibel of noise figure, even if the VGA's own noise figure worsens by
several decibels; reducing it at the LNA would cost far more noise figure. That is the logic of a
gain table.
\end{worked}

\begin{tryit}[Build the line-up yourself]
Lab 36 (\lab{lab36\_rf\_transceiver.py}), experiment \emph{Receiver line-up builder}, reproduces this
table live. Start from the defaults, then drag \emph{Switch + balun loss} from 2.5 to 0.5\,dB and watch
the noise figure fall almost one for one. Now sweep \emph{LNA gain}: the NF curve flattens, the IIP3
curve keeps falling, and the SFDR curve has a peak. Change \emph{Channel bandwidth} from 1\,MHz to
20\,MHz and see the whole SFDR picture slide.
\end{tryit}

\mdcfig{ch07_two_tone}{Left: a two-tone test through a cubic nonlinearity produces
third-order products at $2f_1-f_2$ and $2f_2-f_1$. Right: fundamental and IM3 output against input
power; their extrapolated intersection is the third-order intercept.}

\subsection{Nonlinearity: compression, intermodulation and intercepts}

Turn up a cheap guitar amplifier and the sound gets louder, then harsher, then it just
\emph{crunches}: the amplifier can no longer swing its output far enough, so it flattens the peaks.
Radio amplifiers do exactly the same, and the crunch has two consequences. The signal stops growing
(\emph{compression}), and new frequencies appear that were not there before (\emph{intermodulation}).

Linearity is usually characterised by expanding the amplifier's transfer function as a power series
(Chapter~\ref{ch:signals}), $y=a_1x+a_2x^2+a_3x^3+\cdots$. Driving it with a single tone
$A\cos\omega t$ gives a fundamental of amplitude $a_1A+\tfrac34a_3A^3$; for a compressive device
($a_3/a_1<0$), the gain falls as $A$ grows. The input at which it has fallen by 1\,dB is the
\term{1\,dB compression point} (P1dB), at
\begin{equation}
A_{\text{1dB}}=\sqrt{0.145\,|a_1/a_3|}.
\end{equation}
Driving it with two tones at $\omega_1$ and $\omega_2$ produces, among many products,
\term{third-order intermodulation} (IM3) at $2\omega_1-\omega_2$ and $2\omega_2-\omega_1$ with amplitude
$\tfrac34a_3A^3$. These are the dangerous ones because they fall close to the tones---in the band, where no
filter can remove them (Figure~\ref{fig:ch07_two_tone}). As the input rises by 1\,dB the IM3 products
rise by 3\,dB; the extrapolated input level at which they would equal the fundamental is the
\term{third-order intercept point} (IIP3), at
\begin{equation}
A_{\text{IIP3}}=\sqrt{\tfrac43\,|a_1/a_3|},
\end{equation}
so for a device described by a cubic, IIP3 lies $10\log_{10}(\tfrac43/0.145)=9.6$\,dB above P1dB. Real
amplifiers deviate (higher-order terms matter near compression), but ``IIP3 $\approx$ P1dB + 10\,dB'' is
the engineer's rule of thumb. The intercept is a fiction---no device can actually be driven there---but a
useful one, because it predicts the IM3 level at any lower drive. If each of two tones is at $P_{\text{in}}$
dBm, the IM3 products at the output, referred to the input, are at
\begin{equation}
P_{\text{IM3}}=3P_{\text{in}}-2\,\text{IIP3}\quad\text{dBm}.
\label{eq:ch07:im3}
\end{equation}
In words: every decibel by which you back off the interferers buys you \emph{three} decibels
less intermodulation, and every decibel of IIP3 you design in buys two.

\begin{deeper}[Where 0.145, 4/3 and 9.6\,dB come from]
With $x=A\cos\omega t$, $\cos^3\theta=\tfrac34\cos\theta+\tfrac14\cos3\theta$, so the fundamental is
$(a_1+\tfrac34a_3A^2)A$. Its gain relative to $a_1$ has fallen by 1\,dB when
$1+\tfrac34(a_3/a_1)A^2=10^{-1/20}=0.891$, i.e.\ $A^2=0.145|a_1/a_3|$. With two equal tones,
$x=A(\cos\omega_1t+\cos\omega_2t)$, expanding $x^3$ gives a term $\tfrac34a_3A^3\cos(2\omega_1-\omega_2)t$
(and its twin), while the linear output is $a_1A$. Setting the two equal gives $A^2=\tfrac43|a_1/a_3|$,
which is $10\log_{10}(1.333/0.145)=9.6$\,dB above the compression point. Because the IM3 amplitude
grows as $A^3$, its power in decibels has slope 3, and the two lines of Figure~\ref{fig:ch07_two_tone}
meet at the intercept; equation~\eqref{eq:ch07:im3} is just the equation of the slope-3 line through
that point.
\end{deeper}

\begin{tryit}[Run a two-tone test]
In Lab 36, experiment \emph{Two-tone test and IIP3}, set \emph{Amplifier model} to \emph{Cubic
(textbook)} and drag \emph{Input power per tone} upward in 1\,dB steps: the IM3 products climb three
times as fast as the tones. Then switch to \emph{Soft limiter (real)} and repeat near compression---the
neat slope of 3 bends, which is why real data sheets measure IIP3 well below P1dB. The lab's challenge
asks you to find P1dB to within 0.3\,dB.
\end{tryit}

The second-order term creates products at $\omega_1\pm\omega_2$ and at DC, and the corresponding
\term{IIP2} is defined the same way with slope 2: $P_{\text{IM2}}=2P_{\text{in}}-\text{IIP2}$. In a
superhet, second-order products generally fall far from the IF and are filtered; in a zero-IF receiver
the difference-frequency product of a blocker with itself lands at DC (Section~\ref{sec:ch07:zeroif}),
which is why IIP2 appears on every RFIC data sheet.

\mdcfig{ch07_cascade_tradeoff}{The line-up trade-off. A chain of band filter ($-1.5$\,dB), LNA
(NF 1.2\,dB, IIP3 0\,dBm, gain varied), mixer (gain 6\,dB, NF 10\,dB, IIP3 $+10$\,dBm) and baseband
(gain 30\,dB, NF 15\,dB, IIP3 $+10$\,dBm). Left: cascade noise figure and IIP3. Right: SFDR in 1\,MHz,
which peaks at a modest LNA gain---more gain only trades dynamic range at the top for a vanishing
improvement at the bottom.}

\paragraph{Cascaded intercept.} Two stages in cascade, each weakly nonlinear, have an overall IIP3 given
to a good approximation (assuming the IM3 products add in phase, the worst case) by
\begin{equation}
\frac{1}{\text{IIP3}_{\text{tot}}}\approx\frac{1}{\text{IIP3}_1}+\frac{G_1}{\text{IIP3}_2}+\frac{G_1G_2}{\text{IIP3}_3}+\cdots
\label{eq:ch07:iip3cascade}
\end{equation}
in linear power units. The proof is short: referred to the input of the cascade, the second stage's
intercept is reduced by the gain in front of it, and the IM3 \emph{amplitudes} add. Notice the mirror
image of Friis: there, gain in front \emph{divides} a stage's contribution; here, it \emph{multiplies}
it. Noise wants gain early; linearity wants it late. Figure~\ref{fig:ch07_cascade_tradeoff} shows the
negotiation for a simple chain whose LNA gain is varied: noise figure falls with LNA gain and
saturates; IIP3 falls steadily; the spurious-free dynamic range defined next has a maximum at a moderate
LNA gain.

\subsection{Spurious-free dynamic range}

The \term{spurious-free dynamic range} combines both ends: the range of input powers over which a signal
is above the noise floor while two equal interferers' IM3 products are still below it. Think of it as
the height of the window through which the receiver sees the world (Figure~\ref{fig:ch07_sfdr_window}):
the sill is the noise floor, the lintel is the strongest pair of interferers that do not yet paint
spurious products on the glass. Setting $P_{\text{IM3}}$ of equation~\eqref{eq:ch07:im3} equal to the
noise floor $P_{\text{noise}}$ gives the largest tolerable interferer,
$P_{\text{max}}=(2\,\text{IIP3}+P_{\text{noise}})/3$, and $\text{SFDR}=P_{\text{max}}-P_{\text{noise}}$:
\begin{equation}
\text{SFDR}=\tfrac23\big(\text{IIP3}-P_{\text{noise}}\big)\ \text{dB},
\qquad P_{\text{noise}}=-174+10\log_{10}B+\text{NF}\ \text{dBm}.
\label{eq:ch07:sfdr}
\end{equation}

\mdcside[0.42]{ch07_sfdr_window}{The dynamic-range window of the worked line-up: noise floor
$-109.6$\,dBm in 1\,MHz, IIP3 $-19$\,dBm. Two interferers up to about $-49$\,dBm leave their IM3
below the noise; the window is 60\,dB tall. The intercept itself lies far above anything the receiver
can actually handle.}

A receiver with NF 6\,dB, IIP3 $-10$\,dBm and a 200\,kHz channel has a noise floor of $-115$\,dBm and
an SFDR of 70\,dB. Because the noise floor rises with bandwidth, SFDR always comes with a bandwidth;
quoting it ``in 1\,Hz'' gives large, flattering numbers that are of little use to a system designer.
Read the equation back in words: SFDR improves by only two-thirds of a decibel for every decibel of
IIP3 or noise figure you win. Dynamic range is expensive, which is why the best receivers in the world
(the laboratory spectrum analysers and military surveillance receivers of
Table~\ref{tab:ch07:arch}) are big, power-hungry superheterodynes rather than single chips. It is also
why the factor $\tfrac23$ is worth remembering: a spec sheet that improves NF by 3\,dB has bought you
only 2\,dB of usable window.

\subsection{Blocking and desensitisation}

A strong \emph{blocker} near the wanted channel does harm in four ways even if no intermodulation
product lands on the signal. It \emph{compresses} the LNA or mixer, reducing gain for the weak signal:
a small signal riding on a large one sees the \emph{incremental} gain, which falls faster than the
large-signal gain (Figure~\ref{fig:ch07_desense}); in the simple cubic model the weak signal's gain is
already about 2\,dB down when the blocker reaches P1dB, and it drops steeply beyond. Its presence
forces the \emph{AGC} to reduce gain, raising the effective noise figure. It \emph{mixes with the LO's
phase noise} to spread noise onto the wanted channel (reciprocal mixing, Section~\ref{sec:ch07:recip}).
And in a zero-IF receiver its \emph{IM2} lands at DC.

\mdcside[0.42]{ch07_desense}{Desensitisation in the cubic model $y=a_1x+a_3x^3$. As a blocker
approaches the compression point, the gain seen by a weak signal riding on it (red) falls about
twice as fast as the blocker's own gain (navy): about 2\,dB down when the blocker is at P1dB.}

Cellular standards therefore specify blocking tests in which the receiver must maintain its
sensitivity, within a few decibels, with a specified interferer present, and those tests, rather than
the plain sensitivity test, usually set the design. For GSM, 3GPP TS~45.005 requires a mobile to decode a
wanted signal 3\,dB above reference sensitivity in the presence of an in-band GMSK blocker at levels of
approximately $-43$\,dBm at 600\,kHz to 1.6\,MHz offset, $-33$\,dBm at 1.6 to 3\,MHz and $-23$\,dBm beyond
3\,MHz---a blocker up to roughly 80\,dB stronger than the wanted signal. LTE and NR specify analogous
in-band, out-of-band and narrowband blocking tests in TS~36.101 and TS~38.101-1.

\subsection{Dynamic-range planning and AGC}
\label{sec:ch07:agc}

A \emph{level diagram} (Figure~\ref{fig:ch07_level_diagram}) traces the wanted signal, the noise and the
largest blocker from the antenna to the ADC, stage by stage, and checks three things at every node:
that the noise stays well above the next stage's own noise (so the noise figure is set early), that the
blocker stays below every stage's compression point with margin for its peak-to-average ratio, and that
at the ADC the blocker fits under full scale while the thermal noise sits comfortably above the
quantisation floor. The analog channel filter, which in a zero-IF receiver is the baseband low-pass
filter, attenuates the blocker before the gain that follows it, and that attenuation is what makes the
plan close.

\mdcfig{ch07_level_diagram}{A level diagram for a narrowband zero-IF receiver (200\,kHz channel, cascade
NF about 5\,dB) facing a $-23$\,dBm blocker 3\,MHz away. The baseband filter attenuates the blocker by
22\,dB and the VGA stage's filtering by another 20\,dB, so it arrives about 11\,dB below the ADC's full scale
(an illustrative $+4$\,dBm) while the thermal noise sits more than 30\,dB above an ideal 12-bit
converter's quantisation noise in the channel. Note the mixer, which sees the blocker at about $-7$\,dBm:
its compression point must be well above that.}

\begin{worked}[How many ADC bits does a GSM blocker need?]
Suppose a zero-IF receiver with a 61.44\,MS/s I/Q ADC pair must handle the GSM case above: wanted
signal at $-99$\,dBm needing about 9\,dB SNR in 200\,kHz, and a $-23$\,dBm blocker at 3\,MHz, with
\emph{no} analog filtering of the blocker before the ADC.

The total noise in 200\,kHz, referred to the antenna, must be at most $-99-9=-108$\,dBm; let us demand that
the ADC's contribution be 10\,dB below that, $-118$\,dBm. Full scale must accommodate the blocker plus a
few decibels of headroom for the wanted signal, noise and AGC error, say $-20$\,dBm referred to the
antenna. The required dynamic range \emph{in the channel} is therefore $-20-(-118)=98$\,dB.

The ADC's quantisation noise is spread over the whole complex band of 61.44\,MHz, of which the channel
filter keeps 200\,kHz: a processing gain of $10\log_{10}(61.44/0.2)=24.9$\,dB. The ADC therefore needs a
full-band SNR of $98-24.9=73$\,dB, which by $\text{SNR}=6.02N+1.76$ (Chapter~\ref{ch:pcm}) is
$N\approx11.8$ ideal bits. A real 12-bit converter has an effective number of bits closer to 10 or 11,
so it falls 5 to 10\,dB short.

The cure is analog filtering. A third-order Butterworth baseband filter with a 1\,MHz corner attenuates
a 3\,MHz blocker by about $60\log_{10}3=29$\,dB, which reduces the needed dynamic range to about 70\,dB and
the needed ADC resolution to under 8 effective bits. This is the central bargain of SDR design: every
decibel of analog selectivity before the converter is a decibel the ADC does not have to provide, and a
wideband SDR that skips the analog filter to keep its bandwidth flexible pays for it in bits.
\end{worked}

\mdcfig{ch07_adc_ladder}{The ADC worked example as a staircase. The channel needs 98\,dB; the
processing gain of filtering 61.44\,MHz down to 200\,kHz returns 24.9\,dB, leaving 73\,dB (11.9 ideal
bits) for the converter. A modest third-order analog filter takes off another 29\,dB and lets a
7-bit-effective converter do the job.}

\paragraph{Automatic gain control.} The level plan holds for one input level; the AGC must keep it valid
as the input varies over 100\,dB. Think of it as the iris of an eye, closing in sunlight and opening in
the dark. It measures power at one or more points (after the LNA, after the
baseband filter, at the ADC output) and adjusts gain in steps, with the rule that gain is removed
\emph{as early in the chain as possible} when signals are strong (to protect linearity) and added as
early as possible when they are weak (to protect noise figure). Its time constants are a compromise: a
fast attack prevents clipping when a strong burst arrives; a slow decay avoids gain pumping on
amplitude-modulated signals; and gain changes in the middle of a packet corrupt it, so packet radios
such as Wi-Fi settle the gain during the preamble's short training field and freeze it for the rest of
the packet. The AD9364 offers both a fast-attack AGC suited to bursty TDD signals and a slow AGC for
continuous signals, plus manual gain control, which is what UHD uses by default on the B200.

\begin{pitfall}[AGC in an SDR experiment]
Many SDR experiments go wrong because the gain is wrong, not because the DSP is. Too little gain and
the ADC's noise dominates; too much and a strong signal elsewhere in the captured band clips the ADC, at
which point every signal in the capture is corrupted by clipping products. On a B200, look at the
peak sample magnitude (UHD scales the converter's full scale to $\pm1.0$ in its floating-point
\texttt{fc32} format, so I or Q values near $\pm1$ mean clipping), not only at the signal of interest: the wideband spectrum view of
\lab{gr01\_spectrum\_iq\_capture.py} shows the signals you did not mean to capture but which are using up
your dynamic range.
\end{pitfall}

\mdcphoto[0.62]{ch07_specan}{A bench spectrum analyser showing a single 868\,MHz remote-control
carrier over a 100\,MHz span. A spectrum analyser is the RF engineer's stethoscope: every quantity in
this section---noise floor, compression, intermodulation products, spurs---is measured on a screen
like this one, and an SDR with an FFT display is a (less calibrated) version of the same instrument.}

%======================================================================
\section{Frequency generation}
\label{sec:ch07:synth}
%======================================================================

Every radio is, at heart, a very precise clock with an antenna attached. The carrier frequency, the
sample clock of the converters, the symbol timing of the modem and the timestamps on every sample
all come from one small component, usually smaller than a fingernail and costing less than a coffee.
This section follows the frequency from that crystal to the LO at the mixer, and asks what happens
when the clock is not quite perfect.

\subsection{References: from crystal to GPSDO}

Every frequency in a radio derives from a reference oscillator, almost always a quartz crystal: a slab
of piezoelectric quartz whose mechanical resonance, with unloaded $Q$ of $10^4$ to $10^6$, controls an
oscillator circuit (Figure~\ref{fig:ch07_xo_teardown}). A $Q$ that high means the crystal ``rings'' for
tens of thousands of cycles after a tap, like a fine wine glass rather than a cardboard cup, and that is
what makes its frequency so stable. Its frequency depends on the cut of the crystal, its temperature,
its age and its drive level (Figure~\ref{fig:ch07_crystal_tempco}). The grades are:

\mdcpairany{ch07_xo_teardown}{Inside a packaged crystal oscillator: the round quartz blank (left), held
on two spring contacts, and the small oscillator-and-divider chip that keeps it ringing.}{ch07_crystal_tempco}{A
crystal's frequency against temperature (illustrative AT-cut curve). A TCXO bends the curve flat
electronically; an OCXO parks the crystal at a turning point, where the slope is zero.}

\begin{description}
\item[XO] (plain crystal oscillator): about $\pm10$ to $\pm50$\,ppm over temperature and life. Fine for
a microcontroller; inadequate for most radios without frequency correction from the received signal.
\item[TCXO] (temperature-compensated): a thermistor network or a digital lookup table bends the frequency
back against the crystal's cubic temperature curve, giving about $\pm0.5$ to $\pm2$\,ppm. Every phone and
most SDRs, including the B200, use one.
\item[OCXO] (oven-controlled): the crystal is held at the turning point of its temperature curve, around
70 to 85\,$^\circ$C, in a small oven, giving parts per billion of stability at the cost of watts of heater
power and minutes of warm-up. Base stations and instruments use them.
\item[GPSDO] (GNSS-disciplined oscillator): a TCXO or OCXO steered by a slow PLL to the timing signals of
a GNSS receiver (Chapter~\ref{ch:spreadspectrum}). Over a day its average frequency is accurate to
parts in $10^{12}$ and it provides a pulse-per-second (PPS) time reference; in the short term it is only
as good as its oscillator. Rubidium and caesium standards serve where GNSS cannot be relied on.
\end{description}

A fractional frequency error $\epsilon$ at carrier $f_c$ is a carrier frequency offset of $\epsilon f_c$
(Figure~\ref{fig:ch07_freq_accuracy}). It is also a sample-clock error of $\epsilon f_s$, because the
same reference clocks the converters, and a timing drift of $\epsilon$ seconds per second. One part per
million sounds tiny, but it is about 86\,ms a day: a crystal wristwatch that good would be a fine
watch, and a radio that bad would be useless at millimetre-wave frequencies without help from the
receiver's synchronizers (Chapter~\ref{ch:sync}).

\mdcfig[0.82]{ch07_freq_accuracy}{Worst-case carrier frequency offset produced by references of
different grades, compared with OFDM subcarrier spacings. A TCXO is good enough for Wi-Fi without
correction of more than a fraction of a subcarrier; for LTE and NR the receiver must estimate and remove
the offset (Chapter~\ref{ch:sync}).}

\begin{worked}[Two parts per million at 2.4\,GHz]
A B200's TCXO is specified at about $\pm2$\,ppm. At 2.4\,GHz the worst-case offset is
$2\times10^{-6}\times2.4\times10^9=4.8$\,kHz; two such radios talking to each other can be off by twice
that, 9.6\,kHz. Against a Wi-Fi subcarrier spacing of 312.5\,kHz this is 3\%: a small fraction, which
Wi-Fi's preamble-based CFO estimator removes easily (802.11 allows each station $\pm20$\,ppm, or up to
96\,kHz between two stations at 2.4\,GHz, and the estimator of Chapter~\ref{ch:ofdm} covers it). Against
an LTE subcarrier spacing of 15\,kHz it is 64\%, which would destroy orthogonality: a phone must
estimate the offset from the synchronisation signals before it can demodulate anything, and 3GPP
requires the UE to transmit within $\pm0.1$\,ppm of the carrier it receives (TS~38.101-1), which it
achieves by locking its oscillator to the base station. The sample clock is off by the same 2\,ppm: at
30.72\,MS/s that is 61 samples per second of drift, so a receiver that does not track timing slips a
sample every 16\,ms. For coherent multi-radio experiments, timing measurements and anything that must
agree with an external clock, feed the B200 a 10\,MHz reference and PPS, or fit its GPSDO module.
\end{worked}

\begin{history}[Quartz takes control]
Walter Cady of Wesleyan University built the first quartz-crystal-controlled oscillator in 1921, and
George Washington Pierce of Harvard published the simple one-transistor (originally one-valve)
oscillator circuit that still bears his name in 1923. Warren Marrison and J.~W.~Horton at Bell
Laboratories built the first quartz clock in 1927, and by the 1930s crystal control had replaced the
free-running tuned circuit in broadcast and military transmitters. The Second World War turned crystal
manufacture into a mass industry: US factories produced tens of millions of crystal units for military
radios, and the techniques for cutting, lapping and finishing quartz to a frequency developed then are
recognisably those used today. The temperature-compensated oscillator was the enabler of the mobile
phone; the GPS-disciplined oscillator, from the 1990s, of the synchronised cellular network.
\end{history}

\mdcphoto[0.55]{ch07_crystals}{A manufacturer's display of quartz crystal packages (cropped): leaded
cans of the HC-49 family at the top, surface-mount ceramic packages shrinking towards the bottom
right, and complete oscillators, including a $5\times7$\,mm TCXO, in the lowest row. A radio's whole
frequency plan hangs on one of these.}

\subsection{PLL synthesisers: integer-$N$ and fractional-$N$}

A 40\,MHz crystal is wonderfully stable, but we need 2.437\,GHz---and next second, perhaps,
915.2\,MHz. The \term{phase-locked loop} (PLL) synthesiser is the gearbox that turns one into the other
(Figure~\ref{fig:ch07_fracn}). A voltage-controlled oscillator's output is divided by $N$ and compared in
phase with the reference (divided by $R$) in a phase-frequency detector; a charge pump and loop filter
steer the VCO until the divided output tracks the reference, so
$f_{\text{VCO}}=Nf_{\text{ref}}/R=Nf_{\text{PFD}}$. Chapter~\ref{ch:sync} develops the loop dynamics;
here we care about frequency resolution, spurs and noise.

An \emph{integer-$N$} synthesiser can only step in multiples of $f_{\text{PFD}}$. A GSM synthesiser with
a 200\,kHz channel raster needs $f_{\text{PFD}}=200$\,kHz, and a loop bandwidth necessarily well below
that---perhaps 10 to 20\,kHz---which makes it slow to settle and leaves the VCO's noise uncorrected over
most of the offsets that matter. Worse, the in-band noise of the phase detector and dividers is
multiplied by $N^2$, and $N$ is large when $f_{\text{PFD}}$ is small.

A \term{fractional-$N$ synthesiser} runs the phase detector at a high frequency (tens of megahertz) and
achieves fine resolution by changing the divide ratio dynamically: dividing by $N$ for some reference
cycles and by $N+1$ for others gives an average of $N+\alpha$ for any fraction $\alpha$
(Figure~\ref{fig:ch07_fracn_dither}). Done periodically, this produces large \emph{fractional spurs} at
offsets of $\alpha f_{\text{PFD}}$; done with a \emph{sigma-delta modulator} (Chapter~\ref{ch:pcm})
driving the divider, the instantaneous ratio wanders over $N-3\ldots N+4$ (for a third-order modulator)
in such a way that the quantisation error is pushed to high offset frequencies, where the PLL's low-pass
loop filters it out. A 24-bit modulator with a 40\,MHz PFD has a resolution of
$40\,\text{MHz}/2^{24}\approx2.4$\,Hz. The price is residual shaped noise near the loop bandwidth,
sensitivity of that noise to charge-pump nonlinearity (which folds the shaped noise back in-band), and
``integer-boundary spurs'' when the fraction $\alpha$ is close to 0 or 1. The AD9364 contains two
fractional-$N$ synthesisers with integrated VCOs, one for receive and one for transmit, both locked to
the board's reference.

\mdcfig{ch07_fracn_dither}{How a fractional-$N$ divider counts to 60.3. Top: switching between 60 and
61 in a regular pattern has the right average (red) but its periodicity creates spurs. Bottom: a
MASH 1-1-1 sigma-delta modulator wanders between 57 and 64 in a busy, noise-like pattern with exactly
the same average, so the error becomes high-frequency noise that the loop filter removes.}

\begin{figure}[htbp]\centering
\begin{tikzpicture}[node distance=0.55cm and 0.55cm,
  blk/.style={draw=navy,fill=navy!6,thick,minimum height=0.85cm,minimum width=1.3cm,align=center,font=\sffamily\scriptsize},
  arr/.style={-{Stealth[length=1.8mm]},thick,navy}]
\node[blk,draw=practc,fill=practc!8] (ref) {TCXO\\40 MHz};
\node[blk,right=of ref] (r) {$\div R$};
\node[blk,right=of r] (pfd) {PFD +\\charge\\pump};
\node[blk,right=of pfd] (lf) {loop\\filter};
\node[blk,right=of lf,fill=accent!6,draw=accent] (vco) {VCO\\$2f_{\text{LO}}$};
\node[blk,right=of vco] (q) {$\div2$\\quadrature};
\node[right=0.35cm of q,font=\sffamily\scriptsize,align=left] (out) {LO 0$^\circ$\\LO 90$^\circ$};
\node[blk,below=0.7cm of lf,minimum width=3.2cm] (div) {multi-modulus divider $\div N[k]$};
\node[blk,below=0.7cm of div,minimum width=3.2cm,draw=accent,fill=accent!6] (sd) {$\Sigma\Delta$ modulator (MASH 1-1-1)};
\node[left=0.5cm of sd,font=\sffamily\scriptsize,align=right] (frac) {fractional word\\$\alpha$ (e.g.\ 24 bits)};
\draw[arr] (ref)--(r); \draw[arr] (r)--node[above,font=\tiny]{$f_{\text{PFD}}$}(pfd); \draw[arr] (pfd)--(lf);
\draw[arr] (lf)--node[above,font=\tiny]{$v_{\text{tune}}$}(vco); \draw[arr] (vco)--(q); \draw[arr] (q)--(out);
\draw[arr] ($(vco.east)+(0.25,0)$) |- (div.east);
\draw[arr] (div.west) -| (pfd.south);
\draw[arr] (sd.north) -- node[right,font=\tiny]{$N+\Delta N[k]$, $\Delta N\in\{-3..4\}$}(div.south);
\draw[arr] (frac)--(sd);
\end{tikzpicture}
\caption{A fractional-$N$ synthesiser with sigma-delta control and quadrature LO generation. The
divide ratio is dithered by a sigma-delta modulator so that its average is $N+\alpha$ while its
quantisation noise is shaped to high offsets, where the loop filter removes it. The VCO runs at twice the
LO frequency and a divide-by-two produces accurate $0^\circ$ and $90^\circ$ phases (and keeps the VCO
away from the PA's output frequency). Compare the charge-pump PLL of Chapter~\ref{ch:sync}.}
\label{fig:ch07_fracn}
\end{figure}

\subsection{Oscillator phase noise and Leeson's model}

Imagine tuning an old radio with a hand that trembles very slightly. The station is there, but the
pointer never sits perfectly still: it wobbles a hair this way and that, so the tone you hear warbles.
A real oscillator is exactly such a wobbly hand on the tuning knob (Figure~\ref{fig:ch07_wobbly_lo}).
Its phase wanders randomly, $v(t)=A\cos(2\pi f_0t+\phi(t))$, and its spectrum, instead of a line, has
skirts described by the \term{single-sideband phase noise} $\mathcal{L}(f)$ in dBc/Hz at offset $f$
from the carrier: for small $\phi$, $\mathcal L(f)\approx\tfrac12S_\phi(f)$, half the power spectral
density of the phase.

\mdcfig{ch07_wobbly_lo}{Phase noise as a wobbly hand. Left: the phase error of a simulated oscillator
wanders by tens of degrees over thousands of samples (exaggerated for visibility). Right: its spectrum.
An ideal oscillator is a single line; the real one is a line with skirts, and those skirts decide how
dense a constellation it can carry and how close a strong neighbour can be.}

D.~B.~Leeson's 1966 model explains the skirts' characteristic shape from a feedback oscillator with a
resonator of loaded quality factor $Q_L$ and a sustaining amplifier of noise factor $F$ delivering
signal power $P_s$:
\begin{equation}
\mathcal L(f)=10\log_{10}\left[\frac{FkT}{2P_s}\left(1+\Big(\frac{f_0}{2Q_Lf}\Big)^2\right)\left(1+\frac{f_c}{f}\right)\right].
\label{eq:ch07:leeson}
\end{equation}
Read it from the outside in. Far from the carrier the phase noise is flat at $FkT/2P_s$: the
amplifier's white noise, half of which is phase modulation, relative to the signal power. Inside the
resonator's half-bandwidth $f_0/2Q_L$, the loop integrates that white phase noise into a random walk, and
the spectrum rises as $1/f^2$ (20\,dB/decade). Below the amplifier's flicker corner $f_c$, upconverted
$1/f$ noise adds another factor of $1/f$, giving the $1/f^3$ region. Figure~\ref{fig:ch07_leeson} (left)
shows the consequences: a 2.4\,GHz on-chip LC VCO with $Q_L=8$, $F=15$\,dB and $P_s=0$\,dBm reaches about
$-118$\,dBc/Hz at 1\,MHz offset, typical of integrated VCOs; a better resonator ($Q=30$ or 200) is
11 or 28\,dB quieter; a 40\,MHz crystal with $Q$ of tens of thousands is far quieter still. The lessons are
the ones every oscillator designer lives by: maximise resonator $Q$, maximise signal power, minimise
amplifier noise and flicker.

\paragraph{Phase noise in a PLL.} Locking an oscillator in a PLL changes its phase noise
(Figure~\ref{fig:ch07_leeson}, right): the steady crystal takes the noisy VCO by the hand. The loop acts
as a low-pass filter on the reference's phase noise (multiplied by $N$, i.e.\ raised by $20\log_{10}N$\,dB,
because multiplying frequency by $N$ multiplies phase deviations by $N$) and on the noise of the phase
detector, charge pump and dividers, and as a high-pass filter on the VCO's noise. The phase-detector
contribution is commonly specified by a \emph{normalised in-band phase-noise floor}, a figure of merit
around $-220$ to $-230$\,dBc/Hz for modern synthesiser ICs, to which $10\log_{10}f_{\text{PFD}}+20\log_{10}N$
is added. The optimum loop bandwidth is roughly where the multiplied in-band noise crosses the
free-running VCO noise; making it narrower lets the VCO's rising noise through (in this example, a
50\,kHz loop gives $1.3^\circ$ RMS instead of $0.35^\circ$), wider brings up the in-band floor and the
sigma-delta noise.

\mdcfig{ch07_leeson}{Left: Leeson's model for a 2.4\,GHz VCO with three resonator $Q$s and for a
40\,MHz crystal oscillator. Right: a 2.4\,GHz PLL synthesiser ($N=60$ from a 40\,MHz reference, loop
bandwidth about 400\,kHz). Inside the loop bandwidth the output follows the reference and phase-detector
noise multiplied by $N$ ($20\log_{10}60=35.6$\,dB); outside it, the free-running VCO. Integrated from
1\,kHz to 10\,MHz the output has about $0.35^\circ$ RMS phase error.}

\subsection{Integrated phase noise and its effects}

Integrating $\mathcal{L}(f)$ over the offsets that matter gives the RMS phase error,
\begin{equation}
\sigma_\phi^2=2\int_{f_1}^{f_2}10^{\mathcal{L}(f)/10}\,df\quad(\text{rad}^2),
\label{eq:ch07:pnint}
\end{equation}
where the lower limit is set by what the receiver's carrier-tracking loop removes and the upper by the
signal bandwidth. The equivalent timing jitter is $\sigma_t=\sigma_\phi/(2\pi f_0)$: $0.35^\circ$ at
2.4\,GHz is 0.4\,ps---the time light takes to cross a tenth of a millimetre. Ettus quotes the B200's LO
phase noise as roughly $1^\circ$ RMS at 3.5\,GHz and $1.5^\circ$ at 6\,GHz. Phase noise hurts in three ways.
\begin{itemize}
\item \textbf{Constellation rotation.} Residual phase jitter after carrier tracking rotates the
constellation (Figure~\ref{fig:ch07_phase_noise}), contributing an EVM of approximately
$\sigma_\phi$ in radians: $1^\circ$ is 1.7\% ($-35$\,dB). 64-QAM tolerates a couple of degrees RMS;
1024-QAM, whose outer points are separated by less than $2^\circ$ in angle, needs well under $0.5^\circ$.
\item \textbf{OFDM impairment.} Phase noise slower than an OFDM symbol produces a \emph{common phase
error} that rotates all subcarriers together and is removed with pilots; faster phase noise produces
inter-carrier interference that pilots cannot remove (Chapter~\ref{ch:ofdm}). Because a synthesiser's
phase noise grows as $20\log_{10}f_0$, a 28\,GHz LO is about 20\,dB noisier than a 2.8\,GHz one built the
same way. This is why 5G NR at millimetre-wave uses wider subcarriers (120\,kHz, so that the ICI-causing
part of the spectrum is a smaller fraction) and adds phase-tracking reference signals (PT-RS).
\item \textbf{Reciprocal mixing}, which deserves its own subsection.
\end{itemize}

\mdcfig{ch07_phase_noise}{Left: typical oscillator phase-noise profile with its $1/f^3$, $1/f^2$ and
floor regions. Centre and right: 64-QAM with $0.5^\circ$ and $3^\circ$ RMS residual phase jitter;
the outer points smear along arcs.}

\begin{tryit}[Smear a constellation]
In Lab 1, experiment \emph{Oscillator phase noise}, choose 64-QAM and drag the \emph{Oscillator
linewidth} slider up: the outer points stretch into arcs exactly as in
Figure~\ref{fig:ch07_phase_noise}. Toggle \emph{Carrier tracking on} and widen the
\emph{Tracking-loop bandwidth}---the loop removes the slow wander and the arcs shrink, until noise in
the wider loop starts to hurt.
\end{tryit}

\subsection{Reciprocal mixing}
\label{sec:ch07:recip}

A receiver's LO is not a line but a line with skirts, and a mixer translates \emph{everything} at its
input by \emph{every} component of the LO. A strong blocker at offset $\Delta f$ from the wanted signal
is therefore converted by the LO's phase-noise skirt at offset $\Delta f$ onto the wanted channel,
exactly as if the blocker itself had skirts. The noise added to a channel of bandwidth $B$ is
\begin{equation}
P_{\text{RM}}\approx P_{\text{blocker}}+\mathcal{L}(\Delta f)+10\log_{10}B\quad\text{dBm},
\label{eq:ch07:recip}
\end{equation}
assuming $\mathcal L$ is roughly constant across the channel. In words: the blocker's power, times the
fraction of the LO's power that lies at the right offset, times the bandwidth you collect it in. No
filter after the mixer can remove it, because it sits on top of the wanted signal; only a better LO, or
a filter that attenuates the blocker \emph{before} the mixer, helps (Figure~\ref{fig:ch07_reciprocal}).
The same mechanism operates in reverse in a transmitter, whose phase noise spreads its own carrier into
neighbouring channels, and is one reason why base-station synthesisers are so much quieter than handset
ones.

\begin{analogy}[A blurry flashlight]
Look for a faint star next to a bright street lamp through a smeared pair of glasses. The smear is
not in the sky; it is in your glasses, and it spreads the lamp's glare across the patch of sky where
the star should be. A noisy LO is the smeared lens: it takes a bright blocker that is \emph{not} in your
channel and paints a halo of its light onto the faint target. Cleaning the lens (a quieter LO) or
turning off the lamp (filtering the blocker before the mixer) are the only cures; staring harder at the
smeared image after the fact (filtering after the mixer) is not.
\end{analogy}

\mdcfig{ch07_reciprocal}{Reciprocal mixing. Left: after down-conversion, a $-40$\,dBm blocker 1\,MHz away
carries the LO's phase-noise skirt onto a weak wanted signal at 0\,Hz. With a poor LO
($-100$\,dBc/Hz at 1\,MHz) the skirt buries the wanted signal; with the PLL of
Figure~\ref{fig:ch07_leeson} it is still well above the thermal floor at the wanted channel. Right: rise of
the noise floor in a 200\,kHz channel (NF 5\,dB) against blocker power for three LO phase-noise levels.}

\begin{worked}[Reciprocal mixing with a $-100$\,dBc/Hz LO]
A receiver with a 200\,kHz channel and a 5\,dB noise figure has a thermal noise floor of
$-174+53+5=-116$\,dBm. A blocker of $-40$\,dBm sits 1\,MHz away, and the LO's phase noise at 1\,MHz offset is
$-100$\,dBc/Hz---a mediocre free-running ring oscillator, or a badly designed synthesiser. By
equation~\eqref{eq:ch07:recip} the reciprocal-mixing noise is $-40-100+53=-87$\,dBm, 29\,dB above the
thermal floor: the receiver has lost 29\,dB of sensitivity to a blocker that is not even in its channel.
For the reciprocal-mixing noise to be 10\,dB below the thermal floor, adding only 0.4\,dB, we need
$-40+\mathcal L+53\le-126$, i.e.\ $\mathcal L(1\,\text{MHz})\le-139$\,dBc/Hz. Repeat the calculation with
GSM's $-23$\,dBm blocker at 3\,MHz and a thermal floor of about $-116$\,dBm, and the requirement is
$\mathcal L(3\,\text{MHz})\le-156$\,dBc/Hz---the classic, very demanding GSM synthesiser specification,
which is usually relaxed in practice by accepting some desensitisation in the 3\,dB the blocking test
allows.
\end{worked}

\begin{tryit}[Smear a blocker onto your channel]
Lab 36, experiment \emph{Phase noise and reciprocal mixing}: set \emph{Blocker power} to $-40$\,dBm,
\emph{Blocker offset} to 1\,MHz and \emph{L at 1 MHz offset} to $-100$\,dBc/Hz, and read the
noise-floor rise; it should match the worked example. Then drag \emph{L at 1 MHz offset} down until the
rise falls below half a decibel, and try the 30\,dB/decade \emph{Slope} to see why close-in blockers are
so much more dangerous.
\end{tryit}

\subsection{LO distribution, quadrature generation and pulling}

Generating the LO is only half the job; it must be delivered to the mixers with accurate quadrature and
without coupling into anything else.

\mdcfig{ch07_quadrature_divider}{Quadrature from a divider. A VCO at twice the LO frequency (top)
clocks two flip-flops, one on rising and one on falling edges. Each divides by two, and since the
edges are half a VCO period apart, the two outputs are a quarter of an LO period---exactly
$90^\circ$---apart, whatever the frequency.}

\begin{itemize}
\item \textbf{Quadrature.} The cleanest way to obtain $0^\circ$ and $90^\circ$ LO phases is to run the VCO
at twice (or four times) the LO frequency and divide by two with a pair of flip-flops triggered on
opposite edges; the quadrature accuracy then depends on the duty cycle and matching of the divider, and
is typically within a degree. Polyphase RC networks are the alternative when a $2\times$ VCO is
impractical.
\item \textbf{Pulling.} A VCO is an injection-lockable oscillator: a strong signal near its frequency
pulls it. In a transmitter the PA output, many decibels stronger than anything else on the chip and
modulated, is exactly such a signal if the VCO runs at the carrier frequency. Running the VCO at
$2f_c$, $4f_c$ or $\tfrac43f_c$ and dividing keeps the VCO away from the PA's frequency and its harmonics.
\item \textbf{Coherence between channels.} The two receive channels of a B210 (AD9361) share one RX
synthesiser and are therefore phase-coherent, which makes it a simple two-element array. Two separate
B200s, even sharing a 10\,MHz reference and PPS, have synthesisers whose phases after each retune are
unrelated, so phase-coherent multi-radio experiments need a calibration signal after every tune (and
UHD's timed commands to retune all radios at the same instant).
\end{itemize}

%======================================================================
\section{I/Q impairments and their digital correction}
\label{sec:ch07:iq}
%======================================================================

The philosophy of the modern transceiver is: build analog circuits that are good enough, then
measure and correct the rest in DSP. Nowhere is this clearer than in the I/Q impairments of
direct conversion. Two mixers, two filters, two amplifiers and two ADCs are supposed to be identical
twins; in silicon they are merely siblings, and the small family differences leave fingerprints on
every signal.

\subsection{The imbalance model and the image-rejection ratio}

Let the receiver's I branch be ideal and the Q branch have relative gain $g=1+\epsilon$ and phase error
$\phi$ (in practice both branches are imperfect, but only the difference matters). If the ideal complex
baseband signal is $x=x_I+jx_Q$, the receiver delivers $y=x_I+j\,g(\cos\phi\,x_Q+\sin\phi\,x_I)$, the
model in \texttt{commlib.iq\_imbalance}. Substituting $x_I=(x+x^*)/2$ and $x_Q=(x-x^*)/2j$ and collecting
terms,
\begin{equation}
y=\mu x+\nu x^*,\qquad \mu=\frac{1+g\,e^{j\phi}}{2},\qquad \nu=\frac{1-g\,e^{-j\phi}}{2}.
\label{eq:ch07:iqmodel}
\end{equation}
The received signal is the wanted signal scaled by $\mu$ plus its \emph{complex conjugate} scaled by
$\nu$. Conjugation mirrors the spectrum: a component at baseband frequency $+f$ in $x$ appears at $-f$ in
$x^*$ (Figure~\ref{fig:ch07_image_tone}). The \term{image-rejection ratio} is the ratio of the wanted to
the mirrored power,
\begin{equation}
\text{IRR}=\frac{|\mu|^2}{|\nu|^2}=\frac{1+2g\cos\phi+g^2}{1-2g\cos\phi+g^2}
\approx\frac{4}{\epsilon^2+\phi^2},
\label{eq:ch07:irr}
\end{equation}
where the approximation holds for small errors because $1-2g\cos\phi+g^2=(1-g)^2+2g(1-\cos\phi)\approx
\epsilon^2+\phi^2$. Equation~\eqref{eq:ch07:irr} is identical to the sideband suppression of the phasing
SSB modulator in Chapter~\ref{ch:analog}, as it must be: both are the same mathematics. Read it back:
the image power is a quarter of the \emph{sum of squares} of the two errors, so the larger error
dominates and halving both buys 6\,dB.

\mdcside[0.42]{ch07_image_tone}{A single tone at $+f$ and the ghost that I/Q imbalance creates at
$-f$, for three image-rejection ratios. In a zero-IF capture, any signal sitting at $-f$ must live with
that ghost.}

What does the image look like? Feed a pure tone into an imbalanced receiver and a ghost appears at the
mirror frequency, as many decibels down as the IRR. In a single-signal zero-IF receiver the ghost of
each part of the signal lands on another part of the same signal, a self-inflicted distortion. In an SDR
capture of a busy band, the ghost of a strong station lands wherever its mirror happens to be---perhaps on
the faint signal you were looking for. Ghosts like these are the first thing to suspect when a capture
shows a weak copy of a strong signal on the opposite side of the centre frequency: retune by a few
hundred kilohertz, and a real signal stays put in absolute frequency while the ghost moves the other
way. This section would be short if imbalance were constant, but it drifts with
temperature, gain setting and frequency, which is why every RFIC calibrates it repeatedly.

\begin{worked}[Image rejection from $1^\circ$ and $0.1$\,dB]
A gain imbalance of 0.1\,dB is $g=10^{0.1/20}=1.0116$, $\epsilon=0.0116$; a phase error of $1^\circ$
is $\phi=0.01745$\,rad. The approximation gives $\text{IRR}\approx4/(1.34\times10^{-4}+3.05\times10^{-4})
=9120$, or 39.6\,dB; the exact formula agrees to 0.01\,dB. Notice that the phase error contributes more
than twice as much as the gain error: $1^\circ$ of phase is ``worth'' about 0.15\,dB of gain mismatch. This
is about the best that careful analog design achieves across temperature and frequency without
calibration; 25 to 35\,dB is more typical of an uncalibrated wideband RFIC. A 39.6\,dB IRR limits the
SNR of a mirrored OFDM subcarrier to about 40\,dB---ample for 256-QAM, marginal for the 4096-QAM of Wi-Fi~7,
which therefore always relies on digital correction.
\end{worked}

\mdcfig{ch07_irr}{Image-rejection ratio. Left: contours of IRR against gain and phase imbalance; the
dot marks $0.1$\,dB and $1^\circ$ (39.6\,dB). Right: IRR against phase error for four gain errors. Phase
and gain errors combine as the sum of squares in~\eqref{eq:ch07:irr}, so the larger one dominates.}

What the image does depends on the architecture. In a zero-IF receiver with a single signal centred at
DC, the image is the signal's own mirror: subcarrier $k$ of an OFDM signal receives interference from
subcarrier $-k$, and a single-carrier constellation is skewed into a parallelogram (left of
Figure~\ref{fig:ch07_iq_correction}). In a low-IF receiver, or in an SDR capture containing several
signals, the image of a strong signal at $+f$ falls on whatever weaker signal sits at $-f$, and the
\emph{required} IRR is the power difference plus the SNR needed: a 30\,dB stronger neighbour and a 15\,dB
SNR requirement need 45\,dB of image rejection (Lab~15 measures this in its \emph{Zero-IF vs low-IF}
experiment).

\subsection{Frequency-dependent imbalance}

The model above is \emph{frequency-independent}: the LO's quadrature error and the mixers' gain
mismatch affect all baseband frequencies equally. But the I and Q branches also have separate
baseband filters, TIAs and ADCs whose responses never match perfectly, and at the band edge of a 20 or
56\,MHz-wide filter a tiny difference in corner frequency becomes a significant gain and phase mismatch.
The result is an IRR that varies across the band, worst at the edges
(Figure~\ref{fig:ch07_irr_vs_freq}). Correction then needs a short FIR
filter on one branch (a ``frequency-dependent'' or ``wideband'' I/Q corrector), or, in OFDM receivers,
a per-subcarrier correction estimated from pilots or the preamble, which is natural because the
channel estimator already works per subcarrier.

\mdcfig[0.5]{ch07_irr_vs_freq}{Frequency-dependent imbalance. Fifth-order I and Q baseband filters
whose corners differ by only 1\% (10 and 10.1\,MHz) give a superb IRR near DC but only about 30\,dB near
the band edge; adding a 0.1\,dB, $1^\circ$ LO imbalance caps the IRR at about 40\,dB everywhere.}

\subsection{DC offset: sources and removal}

The DC offset of Section~\ref{sec:ch07:zeroif} has a static part (LO self-mixing, baseband amplifier
offsets) and a dynamic part (time-varying LO reflection, IM2 of a varying blocker). Analog loops remove
the bulk so that the baseband amplifiers do not saturate; a digital estimator removes the rest. The
simplest digital estimator is a running mean subtracted from the signal,
\begin{equation}
m[n]=m[n-1]+\alpha\big(y[n]-m[n-1]\big),\qquad z[n]=y[n]-m[n],
\end{equation}
which is a first-order high-pass filter with a notch at DC whose width is about $\alpha f_s/2\pi$
(Figure~\ref{fig:ch07_dc_notch}). A small $\alpha$ gives a narrow notch that barely touches the signal
but tracks slowly; a large one tracks fast but cuts a hole in the middle of the band. The B200's FPGA
and the AD9364 both implement DC tracking, and UHD enables it by default.

\begin{pitfall}[DC removal eats signals at DC]
Automatic DC-offset correction cannot tell an unwanted offset from a wanted signal that happens to be at
0\,Hz: a carrier tuned exactly to the centre of the band, an unmodulated pilot tone, or the carrier
component of an AM or FSK signal. If your received tone mysteriously vanishes when tuned dead centre,
this is why. The standard SDR answer is \emph{offset tuning}: tune the RF LO a few hundred kilohertz away
from the signal and let the digital down-converter shift it the rest of the way, so that both the DC
spur and the DC notch fall outside the signal (Figure~\ref{fig:ch07_offset_tuning}). In UHD this is a \texttt{tune\_request} with an LO offset;
in GNU Radio, the \texttt{lo\_offset} argument of the USRP source.
\end{pitfall}

\mdcfig[0.5]{ch07_dc_notch}{The running-mean DC remover is a high-pass filter. Each factor of ten in
$\alpha$ moves the notch edge (dotted) by a factor of ten: a narrow notch is gentle on the signal but
slow to follow a changing offset.}

\subsection{Estimating and correcting the imbalance}

Because the imbalance is a $2\times2$ linear transformation of $(x_I,x_Q)$, it can be undone by its
inverse once $g$ and $\phi$ are known. They can be estimated \emph{blindly} from the received signal's
second-order statistics. Almost all communication signals are \emph{proper} (circular): their I and Q
components have equal power and are uncorrelated, $\E[x_I^2]=\E[x_Q^2]$ and $\E[x_Ix_Q]=0$
(Chapter~\ref{ch:noise}). A good constellation looks the same if you rotate it; the imbalance squashes
it into an ellipse, and the squash betrays the error. With $y_I=x_I$ and
$y_Q=g(\cos\phi\,x_Q+\sin\phi\,x_I)$, after removing the mean,
\begin{equation}
\frac{\E[y_Iy_Q]}{\E[y_I^2]}=g\sin\phi,\qquad \frac{\E[y_Q^2]}{\E[y_I^2]}=g^2,
\end{equation}
so $g\sin\phi$ and $g\cos\phi=\sqrt{g^2-(g\sin\phi)^2}$ follow directly, and the corrected quadrature
component is
\begin{equation}
\hat x_Q=\frac{y_Q-g\sin\phi\;y_I}{g\cos\phi}.
\label{eq:ch07:iqcorr}
\end{equation}
This is a Gram--Schmidt orthonormalisation of Q against I, implemented in Lab~1's \texttt{iq\_correct}
function. Figure~\ref{fig:ch07_iq_correction} shows it at work on a capture with 1.5\,dB and $10^\circ$ of
imbalance plus a DC offset; the estimates from 65\,536 samples are correct to 0.01\,dB and $0.01^\circ$. An
equivalent formulation, popular in the literature, estimates the coefficient $w$ of a correction
$z=y+w\,y^*$ that makes $\E[z^2]=0$ (restores circularity), and an adaptive version updates $w$ with an
LMS-style rule, $w\leftarrow w-\mu_wz^2$, tracking slow drift. The method fails for signals that are not
proper---a single real tone, a BPSK signal at DC, a signal whose spectrum is not symmetric in power
within the estimation window---which is why calibration-based methods are also used.

\begin{tryit}[Skew a constellation, then straighten it]
Lab 36, experiment \emph{I/Q imbalance and the image}: with \emph{Signals} set to \emph{Zero-IF: one
16-QAM signal}, drag \emph{Phase error} to $5^\circ$ and watch the square constellation shear into a
parallelogram. Switch to \emph{Capture: strong neighbour at the mirror}, raise \emph{Neighbour above
wanted} to 40\,dB, and then toggle \emph{Blind digital correction}: the ghost collapses by tens of
decibels, exactly the effect of equation~\eqref{eq:ch07:iqcorr}.
\end{tryit}

\mdcfig{ch07_iq_correction}{Blind I/Q correction. A 16-QAM signal at DC and a 15\,dB stronger
QPSK carrier at $+0.27$ cycles/sample are received with 1.5\,dB and $10^\circ$ of imbalance and a DC
offset. Left: the 16-QAM constellation is skewed and shifted; centre: after mean removal and the
correction~\eqref{eq:ch07:iqcorr}; right: the spectrum before and after, showing the DC spike and the
QPSK carrier's image at $-0.27$ removed.}

\paragraph{Calibration.} An RFIC can measure its own impairments by injecting a known test tone, or by
looping its transmitter back into its receiver, at power-up and whenever the frequency or gain changes.
The AD9364 performs a sequence of such calibrations (baseband and RF DC offset, receive and transmit
quadrature) when it is tuned, and runs background tracking of receive quadrature and DC offset during
operation; UHD exposes the tracking as \texttt{set\_auto\_dc\_offset} and \texttt{set\_auto\_iq\_balance}.

\subsection{The transmitter's I/Q impairments}

A quadrature modulator, $s(t)=I(t)\cos2\pi f_ct-Q(t)\sin2\pi f_ct$, suffers the same impairments in
mirror image. A DC offset at the DAC outputs, together with the mixer's LO-to-RF feedthrough, appears as
\term{LO leakage}: a carrier spike at the centre of the transmitted signal. Gain and phase mismatch
produce an image at the mirror frequency, suppressed by the IRR of equation~\eqref{eq:ch07:irr}; the
Ettus data sheet quotes the B200's single-sideband and carrier suppression at roughly 35 and 50\,dBc
respectively. Both are corrected by \emph{pre}-distorting the I/Q samples: adding a small DC term of
opposite sign cancels the leakage, and applying the inverse of~\eqref{eq:ch07:iqmodel} before the DACs
cancels the image. The estimates come from a loopback measurement through the receiver or a dedicated
observation receiver, and both errors count in the transmitter's EVM budget
(Figure~\ref{fig:ch07_lo_leakage}).

\begin{keyidea}[Analog good enough, digital the rest]
Every impairment in this section---DC offset, I/Q imbalance, LO leakage, frequency-dependent
mismatch---is a \emph{deterministic, slowly varying, linear} distortion. That is what makes it cheap to
correct digitally: estimate a handful of parameters occasionally and apply a $2\times2$ matrix or a short
filter to every sample. The analog designer's job is to keep the impairments small enough and stable
enough that the estimates converge and stay valid between updates; the DSP designer's is to estimate
them from whatever the signal offers. The same partnership, with a much harder nonlinear model, is the
basis of the digital predistortion of the next section.
\end{keyidea}

%======================================================================
\mdcfig{ch07_lo_leakage}{A transmitter's I/Q fingerprints, simulated with a single-sideband test tone.
Uncalibrated (red: 0.4\,dB, $3^\circ$ and a small DC offset), the LO leaks through at the centre and an
image appears at the mirror frequency, both about 30\,dB down. After loopback calibration (navy) both
fall by a further 25\,dB or more.}

\section{The transmitter}
\label{sec:ch07:tx}
%======================================================================

The receiver's enemies are noise and strong neighbours; the transmitter's are heat and
\emph{its own} distortion. A phone's power amplifier can account for a large share of the battery
drain during a call or upload, and a macro base station's PAs, with their cooling, are the biggest
item on its electricity bill. Almost everything in this section is a battle to make the PA both linear
(so it does not pollute its neighbours) and efficient (so it does not cook itself)---two goals that
pull in opposite directions.

\subsection{From samples to RF: DACs and reconstruction}

A modern transmitter synthesises the complex envelope digitally, interpolates it to the DAC rate,
converts I and Q with two DACs, filters, and up-converts with the quadrature modulator. A DAC holds each
sample for one sample period, so its output is the ideal impulse train convolved with a rectangular
pulse: the \emph{zero-order hold} of Chapter~\ref{ch:pcm}. Two consequences follow
(Figure~\ref{fig:ch07_dac_images}). First, the output spectrum is multiplied by a sinc envelope, which
droops by $20\log_{10}\sinc(f/f_s)$: 0.6\,dB at $0.2f_s$ and 3.9\,dB at $f_s/2$. An \emph{inverse-sinc}
FIR before the DAC flattens it. Second, the spectrum contains \emph{images} of the signal around every
multiple of $f_s$, attenuated only by the sinc nulls; the image at $f_s-f$ is only a few decibels below the
wanted signal when $f$ is a sizeable fraction of $f_s$. An analog reconstruction filter must remove them.
Interpolating digitally by 4 before the DAC pushes the first image to $4f_s-f$, so that the sinc envelope
attenuates it heavily and a gentle, cheap, low-group-delay analog filter suffices. The AD9364 does
exactly this: its transmit path has a programmable FIR and up to three half-band interpolators before
DACs that run at a multiple of the data rate.

\mdcfig{ch07_dac_images}{DAC images and the zero-order hold. Left: running the DAC at the data rate
leaves images only a few decibels down at $f_s\pm f$, demanding a steep reconstruction filter. Right: after
$4\times$ digital interpolation the first images move to $4f_s\pm f$, deep in the sinc envelope's skirt,
and a gentle analog filter suffices.}

\subsection{Power-amplifier classes}

The PA is the most power-hungry component of most radios and the most nonlinear. Its \emph{drain
efficiency} is $\eta=P_{\text{out}}/P_{\text{DC}}$; its \emph{power-added efficiency} subtracts the input
drive, $(P_{\text{out}}-P_{\text{in}})/P_{\text{DC}}$. The classic classes are distinguished by how the
transistor is biased and driven (Table~\ref{tab:ch07:paclasses}), and most simply by the fraction of
each RF cycle during which it conducts (Figure~\ref{fig:ch07_class_waveforms}). A class-A stage is a
tap always running; a class-C stage opens the tap only for a short squirt at the top of each cycle.

\mdcfig{ch07_class_waveforms}{Conduction angle. The drain current of a class-A stage never stops; a
class-AB stage rests for part of each cycle, class B for half, class C for most of it. The less time
the transistor spends conducting while it also has voltage across it, the less power it burns---and the
less faithfully its output follows the input's amplitude.}

\begin{table}[htbp]\centering\small
\caption{Power-amplifier classes. Efficiencies are ideal maxima at full output.}
\label{tab:ch07:paclasses}
\begin{tabular}{l>{\raggedright\arraybackslash}p{2.5cm}c>{\raggedright\arraybackslash}p{5.6cm}}
\toprule
Class & Operation & $\eta_{\max}$ & Typical use \\
\midrule
A & conducts $360^\circ$ & 50\% & Small-signal drivers, highest-linearity stages \\
AB & conducts $180$--$360^\circ$ & 50--78.5\% & Linear PAs for QAM/OFDM (phones, Wi-Fi, carrier amp of a Doherty) \\
B & conducts $180^\circ$ (push-pull) & 78.5\% & Linear PAs; the theoretical reference \\
C & conducts $<180^\circ$ & $\to100$\% & FM and CW transmitters, peaking amp of a Doherty \\
D & switching, square wave & 100\% & Audio, LF/MF broadcast, envelope modulators \\
E & switching with shaped load (zero-voltage switching) & 100\% & GSM, Bluetooth, RFID, constant-envelope PAs \\
F, F$^{-1}$ & harmonic tuning shapes voltage and current & 100\% & High-efficiency base-station and satellite PAs \\
\bottomrule
\end{tabular}
\end{table}

The efficiency limits follow from simple waveform arguments. A class-A stage biased at $I_{DC}$ with
supply $V_{DD}$ delivers at most a sinusoid of amplitude $V_{DD}$ and $I_{DC}$, i.e.\
$P_{\text{out}}=\tfrac12V_{DD}I_{DC}$, half the DC power. A class-B stage conducts on alternate half
cycles; its DC current is the average of a half-sine, $(2/\pi)I_{\text{pk}}$, so
$\eta=\tfrac12V_{\text{pk}}I_{\text{pk}}/(V_{DD}\cdot2I_{\text{pk}}/\pi)=(\pi/4)(V_{\text{pk}}/V_{DD})$,
reaching 78.5\% at full swing. Switching classes D, E and F make the transistor a switch---either full
voltage and no current, or full current and no voltage---so ideally it dissipates nothing; but a switch
produces a constant-amplitude output, so these classes can amplify only constant-envelope signals, or
must have their amplitude restored by modulating the supply.

\paragraph{Efficiency versus back-off.} The crucial point for modern signals is what happens \emph{below}
full output. If the output amplitude is a fraction $v$ of its maximum (an output back-off of
$-20\log_{10}v$\,dB), the class-A efficiency falls as $v^2$ (the DC power is constant) and the class-B
efficiency as $v$ (the DC current follows the signal). An OFDM signal spends most of its time 6 to 12\,dB
below its peak (Figure~\ref{fig:ch07_pa_efficiency}), so averaged over the signal's amplitude
distribution an ideal class-B PA handling an 8\,dB-PAPR OFDM signal is only about 35\% efficient, and a
class-A one 8\%. Real PAs, with knee voltages, losses and the need for some extra back-off for
linearity, do worse. Constant-envelope modulations (FM, GMSK, FSK) avoid the problem entirely, which is
why 2G GSM, Bluetooth and most IoT links use them: they let a battery-powered device run a saturated,
efficient PA (Chapter~\ref{ch:modulation}).

\mdcfig{ch07_backoff_shout}{Why linear PAs must back off. The envelope of an OFDM signal (navy) has
peaks well above its RMS level. If the PA saturates only 5\,dB above the average, the tallest peaks
(red) are flattened---about 4\% of samples here---and every flattened peak sprays distortion into the
neighbouring channels. Saturating 10\,dB above the average clips almost nothing, but the PA now spends
nearly all its time far below its efficient full-power point.}

\begin{analogy}[Don't shout so loud your voice cracks]
A singer who wants to be heard at the back of the hall sings loudly---but not at the very top of their
range, where the voice cracks on the loudest notes. An OFDM signal is a song with occasional very loud
notes; the PA is the voice. Run it so hard that the peaks hit saturation and the voice cracks (spectral
regrowth, EVM); hold it back so the peaks fit and most of the song is sung at a fraction of full voice,
where the amplifier is inefficient. Crest-factor reduction trims the loudest notes, digital
predistortion trains the voice not to crack, and the Doherty and envelope tracking let a quiet voice be
efficient too. The analogy's limit: a singer's voice gets tired, while a PA's problem is pure heat.
\end{analogy}

\mdcfig{ch07_pa_efficiency}{Left: ideal drain efficiency against output back-off for class A, class B,
symmetric and asymmetric Doherty and envelope tracking (modelled as a 70\%-efficient PA with an
85\%-efficient supply modulator), with the histogram of an 8\,dB-PAPR OFDM signal's instantaneous back-off.
Right: the average efficiency each achieves with that signal.}

\subsection{AM/AM, AM/PM and memory}

A PA's nonlinearity is characterised, for narrowband signals, by how its output amplitude and phase
depend on the input amplitude: the \term{AM/AM} and \term{AM/PM} characteristics. For solid-state PAs the
\emph{Rapp model} captures AM/AM well,
\begin{equation}
|y|=\frac{G|x|}{\big(1+(G|x|/A_{\text{sat}})^{2p}\big)^{1/2p}},
\end{equation}
with a smoothness parameter $p$ of about 2 to 3; travelling-wave tube amplifiers in satellites are
described by the \emph{Saleh model}, which has pronounced AM/PM (Chapter~\ref{ch:satellite}). Both are
\emph{memoryless}: the output depends only on the present input. The smoothness parameter
decides how gently the amplifier runs out of breath (Figure~\ref{fig:ch07_rapp_family}): a small $p$
is a soft knee that distorts a little at every level, a large $p$ an almost perfect limiter that is
linear right up to saturation and then stops dead.

\mdcfig[0.5]{ch07_rapp_family}{The Rapp AM/AM model for four smoothness factors. Solid-state PAs are
typically $p\approx2$--3; $p\to\infty$ is an ideal clipper.}

Real PAs have \term{memory}: the output also depends on past inputs, because the bias networks, matching
networks and the transistor's temperature all have time constants. Electrical memory from bias-network
impedance at the envelope frequency, and thermal memory with time constants from microseconds to
milliseconds, make the AM/AM and AM/PM ``curves'' into clouds and hysteresis loops when measured with a
wideband signal (Figure~\ref{fig:ch07_dpd_amam}, red). Memory effects grow with signal bandwidth, which
is why they matter more for a 100\,MHz NR carrier than for a 5\,MHz LTE carrier.

\subsection{Spectral regrowth, ACLR and EVM}

Linear modulations such as QAM and OFDM have a large peak-to-average power ratio (PAPR), about 10\,dB
for OFDM at $10^{-3}$ probability (Chapter~\ref{ch:ofdm}). A PA driven near saturation compresses the peaks,
and the distortion---odd-order products of the signal with itself---spreads into adjacent channels as
\term{spectral regrowth} (Figure~\ref{fig:ch07_pa_regrowth}). It is the same intermodulation as the
two-tone test of Section~\ref{sec:ch07:metrics}, with thousands of tones intermodulating at once.
Third-order distortion occupies three times the signal bandwidth, fifth-order five times. The
\term{adjacent channel leakage ratio} (Figure~\ref{fig:ch07_aclr_channels}),
\begin{equation}
\text{ACLR}=10\log_{10}\frac{\int_{\text{assigned channel}}S(f)\,df}{\int_{\text{adjacent channel}}S(f)\,df}\ \text{dB},
\end{equation}
is the standard measure. For NR, 3GPP requires about 30\,dB from a phone (TS~38.101-1) and 45\,dB from a
base station (TS~38.104); LTE's numbers (TS~36.101, 36.104) are similar. In-band distortion is measured by
the \term{error-vector magnitude}, the RMS difference between the transmitted and ideal constellation
points relative to the RMS signal; NR allows a phone 8\% for 64-QAM and 3.5\% for 256-QAM. Beyond the
adjacent channel, a \emph{spectrum emission mask} bounds the out-of-band emission as a function of
offset, and \emph{spurious emission} limits (of the order of $-30$\,dBm per megahertz above 1\,GHz for a
phone) apply everywhere else, including at the harmonics that the PA's output filter must suppress.

\mdcfig{ch07_aclr_channels}{What ACLR measures. An OFDM signal after a Rapp PA driven with 6\,dB of
input back-off: the power in the assigned channel (green) is compared with the power that has leaked
into each adjacent channel (red). Here the ratio is about 30\,dB---just the NR handset limit, and
15\,dB short of a base station's.}

\mdcfig{ch07_pa_regrowth}{Left: AM/AM curve of a solid-state PA (Rapp model). Right: spectrum of an
OFDM signal after the PA at 12, 6 and 3\,dB of back-off. Each 3\,dB less back-off raises the
out-of-band emission substantially.}

\mdcfig{ch07_aclr_backoff}{A 64-QAM OFDM signal through a Rapp PA with AM/PM. Left: ACLR rises by about
1.7\,dB per dB of back-off while efficiency (dotted, ideal class B) falls; meeting a base station's 45\,dB
without linearisation needs about 14\,dB of back-off and leaves less than 20\% efficiency. Right: EVM
against back-off with the NR UE limits for 64- and 256-QAM.}

\begin{tryit}[Drive a PA into trouble]
Lab 36, experiment \emph{Power amplifier: ACLR and EVM}: start at an \emph{Input back-off} of 12\,dB
and walk it down to 2\,dB. Watch the shoulders of the spectrum rise, the constellation blur, and the
ACLR and EVM readouts cross the 3GPP limits (\emph{Show 3GPP limits} on). Then raise \emph{AM/PM at
saturation} and lower \emph{PA smoothness (Rapp p)} to see a softer, more ``tube-like'' amplifier
misbehave at gentler drive.
\end{tryit}

\begin{worked}[The efficiency--ACLR trade for a base-station PA]
A macro base station must deliver 40\,W average per carrier with 45\,dB ACLR. Take the PA model of
Figure~\ref{fig:ch07_aclr_backoff} as representative of a class-AB device without linearisation: 45\,dB
ACLR needs about 14\,dB of output back-off from saturation, where the efficiency is about 18\%. The PA
must then be capable of $40\times10^{1.4}\approx1000$\,W peak, and it draws $40/0.18=220$\,W from the
supply, dissipating 180\,W as heat.

Now add DPD that improves ACLR by about 20\,dB (Figure~\ref{fig:ch07_dpd}). The PA can run at about 6\,dB
back-off, where its raw ACLR of 30\,dB becomes about 50\,dB after correction and its efficiency is about
41\%: a 160\,W device drawing 98\,W. The saving is 120\,W per carrier, about 1\,MWh a year, before counting
the air-conditioning that would have removed the heat. Replace the class-AB stage with a Doherty and the
efficiency at that back-off rises to 50 to 60\% in practice. Multiply by three sectors, several carriers
and hundreds of thousands of sites, and DPD plus Doherty is worth whole power stations; this is why every
macro base station since the late 2000s uses both.
\end{worked}

\begin{worked}[An EVM budget for 256-QAM]
EVM contributions from independent impairments add in power. Suppose a transmitter must meet 3.5\% for
256-QAM, and its designer allocates: LO phase noise $1^\circ$ RMS, contributing
$\sigma_\phi=0.0175$, i.e.\ 1.75\%; I/Q image after calibration at 40\,dB IRR, 1.0\%; PA distortion
after DPD, 2.0\%; DAC quantisation, filter ripple and group delay, 0.5\%. The total is
$\sqrt{1.75^2+1.0^2+2.0^2+0.5^2}=2.9$\%, leaving a little margin for production spread and temperature.
The budget shows where effort pays: halving the phase noise to $0.5^\circ$ would save more than any
improvement to the DAC, and the PA, as usual, is the biggest single term.
\end{worked}

\mdcside[0.42]{ch07_evm_budget}{The 256-QAM EVM budget of the worked example. Contributions add as the
root sum of squares, so the total (2.88\%) is far less than their plain sum (5.25\%) and is dominated by
the two largest terms, the PA and the phase noise.}

Notice how the budget behaves. Because the contributions add in \emph{power}, the total is set by the
biggest one or two terms, and shaving a small term achieves almost nothing: cutting the 0.5\% DAC term
to zero improves the total from 2.88\% to 2.84\%. Halving the 2\% PA term, on the other hand, brings it
to about 2.3\%. This is why EVM budgets are negotiated between teams at the start of a project---the
synthesiser designers, the PA designers and the DSP engineers who own the predistorter---and why a
transceiver that misses its EVM target almost always misses it in the PA or the oscillator. The same
root-sum-square logic will reappear when we build link budgets with several independent noise sources
in Chapter~\ref{ch:channels}, and in satellite payloads, where intermodulation in the transponder is
budgeted alongside thermal noise (Chapter~\ref{ch:satellite}). In each case the engineering rule is the
same: find the biggest term, attack it, and stop polishing the small ones.

\subsection{Crest-factor reduction}

\term{Crest-factor reduction} (CFR) reduces the PAPR of the digital signal before the PA, trading a
small, controlled amount of EVM for a large gain in efficiency---trimming the loudest notes of the song
before the singer has to sing them. The simplest effective method is \emph{iterative clipping and
filtering}: clip the complex envelope's magnitude at a threshold, then filter away the out-of-band
products the clipping created (for OFDM, zero the unused subcarriers), and repeat, because filtering
regrows some peaks. Figure~\ref{fig:ch07_cfr} shows a 64-QAM OFDM signal whose PAPR at $10^{-4}$
probability falls from about 10\,dB to under 8\,dB with one iteration and to about 7\,dB with four, at
EVM costs of 3.3\% and 5.7\%. Production CFR uses \emph{peak cancellation}---subtracting
band-limited pulses centred on each peak---which achieves the same result in a single pass and can shape
the EVM, putting more of it on low-order subcarriers that can tolerate it. Tone reservation (Chapter~\ref{ch:ofdm})
keeps EVM at zero by sacrificing a few subcarriers.

\mdcfig{ch07_cfr}{Crest-factor reduction by clipping and filtering. Left: CCDF of the instantaneous-to-average
power for the original OFDM signal and after one and four clip-and-filter iterations at a 6\,dB
threshold. Right: the spectrum per subcarrier; clipping alone (dashed) splatters into the guard band,
while filtering confines the distortion to the occupied subcarriers.}

\begin{tryit}[Trim the peaks]
Lab 36, experiment \emph{Crest-factor reduction}: lower \emph{Clip level above rms} from 9 to 5\,dB
with \emph{Clip-and-filter iterations} at 1, and watch the CCDF pull in. Turn \emph{Filter after
clipping} off to see the splatter of Figure~\ref{fig:ch07_cfr}; turn it back on and raise the
iterations to recover the peaks that the filter regrows.
\end{tryit}

\subsection{Digital predistortion}
\label{sec:ch07:dpd}

Suppose you must sign a cheque with a pen that always smears your downstrokes to the right. A clever
forger would learn the pen's habit and write each stroke slightly to the \emph{left}, so the smear
lands exactly where it should. \term{Digital predistortion} (DPD) is that forger. It applies, in the
digital domain, an approximate inverse of the PA's nonlinearity, so that the cascade of predistorter and
PA is linear up to saturation (Figure~\ref{fig:ch07_dpd_handwriting}).

\mdcfig{ch07_dpd_handwriting}{The idea of predistortion. The PA compresses (centre); the predistorter is
deliberately expansive (left), boosting large amplitudes by exactly what the PA will take away; the
cascade is a straight line (right) up to the point where the PA simply has no more power to give.}

\begin{analogy}[Pre-distorting your handwriting]
DPD is handwriting adjusted for a bad pen: if the pen makes every stroke too short at the end, write
the ends a little longer, and the page looks right. The adjustment must be learned by watching what the
pen actually does---that is the observation receiver---and relearned as the pen warms up and wears.
The analogy's limit is saturation: if the pen has run out of ink, no amount of clever handwriting puts
ink on the page, and no predistorter can make a PA deliver more than its maximum power.
\end{analogy}

Because the PA has memory, the predistorter must too. The workhorse model is the \term{memory polynomial},
\begin{equation}
z[n]=\sum_{k=1}^{K}\sum_{m=0}^{M}a_{km}\,x[n-m]\,\big|x[n-m]\big|^{k-1},
\label{eq:ch07:mp}
\end{equation}
a polynomial of order $K$ in the envelope at each of $M+1$ delays, with complex coefficients $a_{km}$.
It is a pruned Volterra series that keeps only the ``diagonal'' terms, and its great virtue is that it is
\emph{linear in its coefficients}: stacking the basis functions $x[n-m]|x[n-m]|^{k-1}$ as the columns of
a matrix $\boldsymbol\Phi$, the coefficients that best map one signal to another are a least-squares
solution, $\mathbf a=(\boldsymbol\Phi^H\boldsymbol\Phi)^{-1}\boldsymbol\Phi^H\mathbf z$. The
\emph{generalised} memory polynomial of Morgan and colleagues adds cross terms between the signal and
lagging and leading envelope samples, capturing more of the memory at modest cost.

\begin{figure}[htbp]\centering
\begin{tikzpicture}[node distance=0.6cm and 0.75cm,
  blk/.style={draw=navy,fill=navy!6,thick,minimum height=0.9cm,minimum width=1.6cm,align=center,font=\sffamily\scriptsize},
  arr/.style={-{Stealth[length=1.8mm]},thick,navy}]
\node[font=\sffamily\scriptsize] (x) {$x[n]$};
\node[blk,right=0.6cm of x] (pd) {predistorter\\(copy of $\mathbf a$)};
\node[blk,right=of pd] (dac) {DAC +\\up-conv.};
\node[blk,right=of dac,fill=accent!6,draw=accent] (pa) {PA};
\node[right=0.6cm of pa,font=\sffamily\scriptsize] (y) {$y(t)$};
\node[blk,below=0.75cm of pa] (obs) {observation\\receiver, $\div G$};
\node[blk,below=0.75cm of pd] (post) {postdistorter\\training: fit $\mathbf a$};
\draw[arr] (x)--(pd); \draw[arr] (pd)--node[above,font=\tiny]{$z[n]$}(dac); \draw[arr] (dac)--(pa); \draw[arr] (pa)--(y);
\draw[arr] ($(pa.east)+(0.25,0)$) |- (obs.east);
\draw[arr] (obs)--node[above,font=\tiny]{$y[n]/G$}(post);
\draw[arr,accent] ($(pd.east)+(0.3,0)$) -- ++(0,-0.75) -| ($(post.north)+(0.45,0)$);
\node[font=\tiny,text=accent,anchor=west] at ($(pd.east)+(0.35,-0.45)$) {target $z[n]$};
\draw[arr,dashed] ($(post.north)-(0.3,0)$) -- node[left,font=\tiny]{copy $\mathbf a$} ($(pd.south)-(0.3,0)$);
\end{tikzpicture}
\caption{The indirect learning architecture. A postdistorter is trained to map the PA's normalised
output $y/G$ back to the PA's input $z$; since a $p$th-order inverse is (approximately) the same whether
placed before or after the PA, its coefficients are copied into the predistorter. Training repeats
every few milliseconds to seconds to track temperature and ageing.}
\label{fig:ch07_ila}
\end{figure}

How are the coefficients found? The PA's output is observed through an \emph{observation receiver}---a
second, linear down-converter and ADC whose bandwidth must be three to five times the signal bandwidth to
see the distortion products. In the \term{indirect learning architecture} (ILA,
Figure~\ref{fig:ch07_ila}), a \emph{postdistorter} is fitted to map the normalised PA output $y/G$ to the
PA input $z$, by least squares on a block of samples; its coefficients are then copied into the
predistorter, and the process iterates a few times. The ILA relies on the fact that the $p$th-order
inverse of a nonlinear system is the same whether applied before or after it; it is simple and robust,
though slightly biased by noise in the observed output. \emph{Direct learning} instead adjusts the
predistorter to minimise the error between $x$ and $y/G$ directly, which avoids the bias at the cost of
needing a model of the PA's gradient.

Figure~\ref{fig:ch07_dpd} shows ILA at work on a PA model with memory (a Wiener--Hammerstein cascade of a
short filter, a Rapp AM/AM with AM/PM, and another filter) driven by OFDM at 11\,dB input back-off. The PA
alone gives 39\,dB ACLR; a memoryless polynomial predistorter ($K=7$, $M=0$) improves it to 56\,dB; a memory
polynomial ($K=7$, $M=3$) to 60\,dB. Figure~\ref{fig:ch07_dpd_amam} shows why: after DPD the AM/AM and
AM/PM clouds collapse onto the ideal straight lines. DPD cannot help above saturation---no predistorter can
make the PA deliver more power than it has---so in practice it is always combined with CFR, which keeps
the peaks below the point where the PA runs out.

\mdcfig[0.88]{ch07_dpd}{Digital predistortion of a PA with memory at 11\,dB input back-off, using the
indirect learning architecture. A memoryless polynomial removes most of the regrowth; a memory polynomial
removes more, approaching the ideal linear amplifier near the channel.}

\mdcfig{ch07_dpd_amam}{AM/AM and AM/PM of the PA with memory, measured sample by sample with an OFDM
signal. Without DPD (red) the characteristic compresses and memory spreads it into a cloud; with the
memory-polynomial DPD (green) it collapses onto the ideal linear response.}

\begin{tryit}[Teach a predistorter, one click at a time]
Lab 36, experiment \emph{Digital predistortion}: with \emph{PA has memory} on, \emph{Nonlinearity order K}
at 7 and \emph{Memory depth M} at 0, press \emph{Run one ILA iteration} a few times and watch the
regrowth fall. Press \emph{Reset the predistorter}, set $M=3$ and repeat: the extra memory taps buy
several more decibels, as in Figure~\ref{fig:ch07_dpd}. Finally lower \emph{Input back-off} towards
4\,dB and find where DPD stops helping---the PA has run out of ink.
\end{tryit}

\begin{inpractice}[DPD in real products]
DPD began in cellular base stations in the 2000s, where the savings of the worked example above paid for
the observation receiver and the FPGA logic many times over, and it is now universal in macro and small
cells, in cable-TV and satellite amplifiers, and in high-power broadcast transmitters. Commercial
implementations run the predistorter at three to five times the signal bandwidth in an FPGA or ASIC and
the coefficient estimation in software on an embedded processor; RFSoC-based radio units and dedicated
transceiver chips include DPD engines as hard blocks. Phones were long thought too power-constrained for
DPD, but the wide bandwidths and high peak-to-average ratios of NR uplinks have brought lightweight DPD
and envelope-tracking-aware predistortion into high-end handset PA modules. Massive-MIMO arrays with 64 or
more low-power PAs pose a new problem: a full DPD per PA is expensive, and research and products are
exploring shared observation receivers and beam-domain DPD.
\end{inpractice}

\subsection{The Doherty amplifier}

The \term{Doherty amplifier} (Figure~\ref{fig:ch07_doherty_tikz}) combines a \emph{carrier} (main)
amplifier biased in class AB with a \emph{peaking} (auxiliary) amplifier biased in class C that turns on
only for signal peaks. Think of two wrestlers in a tag team: the carrier amplifier does all the work
while the going is easy, and the peaking amplifier leaps into the ring only when the opponent (the
signal peak) gets big---and, crucially, its arrival changes the conditions so that its partner can
work harder too.

\begin{figure}[htbp]\centering
\begin{tikzpicture}[node distance=0.6cm and 0.8cm,
  blk/.style={draw=navy,fill=navy!6,thick,minimum height=0.9cm,minimum width=1.6cm,align=center,font=\sffamily\scriptsize},
  tl/.style={draw=practc,fill=practc!8,thick,minimum height=0.55cm,minimum width=1.5cm,align=center,font=\sffamily\scriptsize},
  arr/.style={-{Stealth[length=1.8mm]},thick,navy}]
\node[font=\sffamily\scriptsize] (in) {RF in};
\node[blk,right=0.5cm of in,minimum width=1.4cm] (split) {splitter};
\node[blk,right=of split,yshift=0.9cm] (main) {carrier amp\\class AB};
\node[tl,right=0.5cm of split,yshift=-0.9cm] (ph) {$\lambda/4$ (90$^\circ$)};
\node[blk,right=0.4cm of ph,fill=accent!6,draw=accent] (peak) {peaking amp\\class C};
\node[tl,right=of main] (inv) {$\lambda/4$, $Z_0=R_{\text{opt}}$\\impedance inverter};
\coordinate (node) at ($(inv.east)+(0.7,0)$);
\node[draw=navy,thick,circle,fill=navy,inner sep=1.3pt] (nd) at (node) {};
\node[blk,right=0.8cm of nd,minimum width=1.4cm] (load) {load\\$R_{\text{opt}}/2$};
\draw[arr] (in)--(split); \draw[arr] (split.east) -- ++(0.2,0) |- (main.west);
\draw[arr] (split.east) -- ++(0.2,0) |- (ph.west);
\draw[arr] (ph)--(peak); \draw[arr] (main)--(inv); \draw[thick,navy] (inv.east)--(nd);
\draw[arr] (peak.east) -| (nd); \draw[arr] (nd)--(load);
\node[above=0.12cm of nd,font=\tiny,align=center] {common node:\\load modulation};
\end{tikzpicture}
\caption{The Doherty amplifier. The quarter-wave line after the carrier amplifier inverts the impedance
at the common node, so that as the peaking amplifier injects current the carrier amplifier sees its load
fall from $2R_{\text{opt}}$ to $R_{\text{opt}}$. The quarter-wave line at the peaking amplifier's input
restores phase alignment at the combining node.}
\label{fig:ch07_doherty_tikz}
\end{figure}

The trick is \emph{load modulation}. The main amplifier drives the load through a
quarter-wave transmission line, an impedance inverter: $Z_{\text{in}}=Z_0^2/Z_L$. At low drive the peaking
amplifier is off, the main amplifier sees twice its optimum load, and so reaches its maximum voltage swing,
and hence maximum efficiency, at half its maximum current---6\,dB below the combined peak power. Above that
point the peaking amplifier starts injecting current into the common node, which raises the apparent load
impedance there; through the inverter the main amplifier sees its load \emph{fall} toward $R_{\text{opt}}$,
so it delivers more current while its voltage stays at the maximum. At full drive both amplifiers deliver
equal power into their optimum loads (Figure~\ref{fig:ch07_doherty}, left). The efficiency of an ideal
symmetric Doherty with class-B devices has two peaks of 78.5\%, at full power and at 6\,dB back-off, and
does not fall below about 70\% between them:
\begin{equation}
\eta(v)=\begin{cases}\dfrac{\pi}{4}\dfrac{v}{\alpha}, & v\le\alpha,\\[2ex]
\dfrac{\pi}{4}\dfrac{v^2}{v(1+\alpha)-\alpha}, & v>\alpha,\end{cases}
\label{eq:ch07:doherty}
\end{equation}
where $v$ is the normalised output amplitude and $\alpha$ the amplitude at which the peaking amplifier
turns on ($\alpha=\tfrac12$ for a symmetric Doherty). An asymmetric Doherty, with a peaking device twice
the size of the main ($\alpha=\tfrac13$), moves the back-off peak to 9.5\,dB, better matched to OFDM's
statistics; three-way Doherties go further. With the 8\,dB OFDM signal of
Figure~\ref{fig:ch07_pa_efficiency}, both ideal Doherties average about 61\%, against 35\% for class B.

\begin{analogy}[Two amplifiers sharing the load like a tag team]
In a tag-team match, one wrestler handles the easy rounds alone; when the opponent gets too big, the
partner jumps in, and the two share the load. The Doherty's twist is that the partner's arrival
\emph{changes the ring}: by pushing current into the shared node it lowers the load the main amplifier
sees, letting the main amplifier deliver more current without running out of voltage swing. Neither
wrestler ever has to fight in an inefficient half-hearted way. The analogy's limit: real tag-team
partners take turns, whereas above the transition both Doherty amplifiers work at once, and making
their hand-over smooth (rather than a lurch in gain and phase) is exactly why Doherty PAs need DPD.
\end{analogy}

\mdcfig{ch07_doherty}{The ideal symmetric Doherty. Left: as the drive rises past half its maximum, the
peaking amplifier's current rises from zero while the main amplifier's load falls from $2R_{\text{opt}}$
to $R_{\text{opt}}$ and its voltage stays at maximum swing. Right: efficiency~\eqref{eq:ch07:doherty} of
symmetric and asymmetric Doherty amplifiers against class B.}

\begin{history}[Doherty's amplifier, 1936]
William H.\ Doherty of Bell Telephone Laboratories described ``a new high efficiency power amplifier for
modulated waves'' in the \emph{Proceedings of the IRE} in 1936. His problem was the AM broadcast
transmitter, whose carrier sits at a quarter of peak envelope power and whose linear amplifiers wasted
most of their input as heat. A Doherty amplifier went into service at a 50\,kW broadcast station, and the
technique was used in AM transmitters for a few decades before high-level plate modulation and later
pulse-step modulation (Chapter~\ref{ch:analog}) displaced it. It then lay largely dormant until the
2000s, when the 3G and 4G base-station industry, facing high-PAPR signals and enormous electricity bills,
rediscovered it. Combined with DPD to correct its rather nonlinear behaviour, it became the standard
base-station PA within a decade. Its contemporary, Henri Chireix's 1935 \emph{outphasing} amplifier,
which combines two constant-envelope signals of varying phase to synthesise an amplitude-modulated one,
has been rediscovered in the same way.
\end{history}

\mdcpairany{ch07_bbc_tx}{A British broadcast transmitter of the early 1920s, used for some of the
first BBC broadcasts (Science Museum stores). Valve transmitters like this turned most of their input
power into heat---the problem Doherty set out to solve.}{ch07_rrh}{A remote radio head for LTE, mounted
next to the antenna on a mast. Its fins dump the heat of its power amplifiers; radio units of this kind
typically combine Doherty PAs, CFR and DPD to keep that heat down.}

\subsection{Envelope tracking and polar transmitters}

\term{Envelope tracking} (ET) attacks the efficiency problem from the supply side: a fast DC--DC
converter (the \emph{supply modulator}) varies the PA's drain voltage to follow the signal's envelope,
so that the transistor operates near compression, where it is efficient, at every instant
(Figure~\ref{fig:ch07_et_supply}). The PA remains a linear amplifier; the supply merely removes the
headroom it would otherwise waste. The supply modulator, which must track envelopes with bandwidths of
tens of megahertz at high efficiency, is the hard part, and the combined efficiency is the product of
the PA and modulator efficiencies (about 58\% in the idealised model of Figure~\ref{fig:ch07_pa_efficiency}).
ET entered smartphones with LTE around 2013 and is now in nearly every 4G and 5G handset's PA module; it
is used with DPD-like ``isogain'' shaping tables that choose the supply voltage for each envelope value
so that the PA's gain stays constant. Its ancestor, \emph{envelope elimination and restoration} (Kahn,
1952), drives a saturated PA with a constant-envelope phase signal and imposes the amplitude entirely
through the supply---the high-power version of the pulse-step AM transmitters of
Chapter~\ref{ch:analog}, and of the polar transmitters used in some EDGE phones.

\mdcfig{ch07_et_supply}{Envelope tracking. With a fixed supply (red line) everything between the
envelope and the supply (pink) is voltage the transistor must drop and turn into heat. A supply that
follows the envelope with a small headroom (green) wastes far less; a minimum voltage keeps the PA
alive in the quiet moments.}

\begin{tryit}[Race the PA architectures]
Lab 36, experiment \emph{PA classes, Doherty and ET}: choose the \emph{OFDM} signal and slide
\emph{PAPR after CFR} from 12 down to 6\,dB, watching the average efficiency of each architecture
climb. Move \emph{Doherty transition $\alpha$} from 0.5 towards 0.33 and see the back-off peak slide to
9.5\,dB; then switch \emph{Signal} to \emph{Constant envelope} to see why GSM never needed any of this.
\end{tryit}

\subsection{Full duplex and self-interference}
\label{sec:ch07:fullduplex}

Every radio in this chapter avoids hearing its own transmitter by separating transmit and receive in
time (TDD) or frequency (FDD). \emph{In-band full duplex}---transmitting and receiving on the same
frequency at the same time---would double spectral efficiency for some links, but it requires suppressing
the transmitter's own signal at the receiver by about 110\,dB: from $+20$\,dBm transmitted to below a
$-90$\,dBm noise floor. That is like hearing a pin drop while you yourself are shouting into the
microphone. No single technique provides this. The suppression is built up in stages
(Figure~\ref{fig:ch07_fullduplex_stack}): antenna isolation or a circulator (20 to 40\,dB), \emph{analog
cancellation} that subtracts a tapped, delayed and weighted copy of the transmit signal at the LNA input
before it can saturate the receiver (20 to 30\,dB), and \emph{digital cancellation} that models the
residual---including the PA's nonlinearity and the transmitter's phase noise, which a linear canceller
cannot remove---with the same memory polynomials as DPD (30 to 40\,dB). Full duplex is used today in
wireline (every DSL and Ethernet transceiver cancels its own echo, Chapter~\ref{ch:wireline}, and
DOCSIS~4.0 extends it to cable), in some relay and integrated-access-backhaul nodes, and in radar and
sensing systems that must listen while transmitting; for general cellular access it remains a research
topic (Chapter~\ref{ch:sixg}).

\mdcfig{ch07_fullduplex_stack}{The full-duplex staircase, taking the upper end of each stage's typical
range. A $+20$\,dBm transmitter must be suppressed by 110\,dB to reach a $-90$\,dBm receiver noise floor:
40\,dB from antenna isolation, 30\,dB from analog cancellation and 40\,dB from digital cancellation.}

%======================================================================
\section{The software-defined radio}
\label{sec:ch07:sdr}
%======================================================================

\subsection{When the radio became software}

Twenty-five years ago a camera was a box of glass, springs and chemistry; today, for most people, a
camera is an app. The lens and sensor are still physical, but everything that makes a photo
\emph{good}---exposure, focus, colour, noise removal, the panorama and the night mode---is software,
and it improves every time your phone updates. Radio has been through the same transformation
(Figure~\ref{fig:ch07_camera_app}).

A \term{software-defined radio} (SDR) is a radio in which some or all of the physical-layer functions
are implemented in software or reprogrammable logic, so that the same hardware can be reconfigured to
different waveforms, bands and standards. In a sense every radio in this chapter is one: the zero-IF
transceiver of Figure~\ref{fig:ch07_transceiver} does everything right of its converters digitally. The
term is usually reserved, though, for radios whose \emph{waveform} is software---where changing from FM
to LTE to a home-made OFDM link is a matter of loading a program, not of buying a chip.

\begin{analogy}[The radio becoming software, like the camera becoming an app]
In a smartphone camera the glass still bends the light and the sensor still counts photons, but what
you \emph{see} is decided by code. In an SDR the antenna, LNA, mixers and converters still do their
physical jobs, but what the radio \emph{is}---FM receiver, GPS receiver, LTE base station, radar---is
decided by the program you run. The analogy's limit is sharp, and this chapter has spent forty pages on
it: a camera app cannot fix a scratched lens, and no SDR program can fix a noise figure, an intercept
point or an oscillator that the hardware does not have.
\end{analogy}

\mdcside[0.42]{ch07_camera_app}{How much of the radio is software? An illustrative estimate of the
share of signal-chain functions done in analog hardware and in digital logic or software, from a 1980s
superhet to a direct-sampling RFSoC radio.}

The story of how radio became software spans seventy years of hardware and three decades of software
(Figure~\ref{fig:ch07_sdr_timeline}). Its first chapter belongs to oscillators, mixers and amplifiers---the
quartz crystal, the Doherty amplifier, Leeson's model, Gilbert's multiplier---which made good analog
radios possible. Its second chapter began in the 1990s, when converters and processors became fast
enough that Joseph Mitola could seriously propose putting the ADC near the antenna and doing everything
else in software. The third chapter is the democratic one: an open-source signal-processing framework,
a university-priced radio board, and finally a \$20 television dongle that put SDR in the pocket of
anyone with a laptop. Each chapter made the next one possible, and each is still visible in the B200:
a crystal, a fractional-$N$ synthesiser, a passive switching mixer, a zero-IF RFIC, an FPGA and an
open-source driver.

\mdcfig{ch07_sdr_timeline}{From quartz to RFSoC: the milestones of this chapter. The analog
inventions on the left made good radios possible; the events on the right (red) turned radios into
software.}

\begin{history}[From SPEAKeasy to the software radio]
The phrase \emph{software radio} was introduced by Joseph Mitola III in the early 1990s (his 1992 IEEE
National Telesystems Conference paper used it) and set out in his influential 1995 \emph{IEEE
Communications Magazine} article ``The software radio architecture'', which described a radio whose ADC
sat as close to the antenna as possible and whose every function, from channel selection to the
protocol stack, was software running on general-purpose hardware. Mitola later coined \emph{cognitive
radio} (1999, with Gerald Maguire) for a software radio that senses and adapts to its environment. The
motivation in the US military was interoperability: the often-told story of the 1983 Grenada operation,
in which units of different services could not talk to each other because their radios were
incompatible, became the standard justification. DARPA and the US Air Force's SPEAKeasy programme
(phase~I from about 1992 to 1995) demonstrated a programmable radio covering 2 to 2000\,MHz that could
emulate a large number of military waveforms. Its successor, the Joint Tactical Radio System (from
1997), and its Software Communications Architecture (SCA), a framework for portable waveform software,
had a long and expensive history, but the SCA is today an international standard for military SDRs, and
the radios fielded under its descendants are software-defined throughout.
\end{history}

\begin{history}[GNU Radio, the USRP and the \$20 dongle]
In civil communications the decisive step came from the open-source world. In 2001 Eric Blossom
started GNU Radio, a free signal-processing framework, with funding from John Gilmore; its early
showpiece was an HDTV receiver. GNU Radio needed hardware, and Matt Ettus built it: the first Universal
Software Radio Peripheral (USRP), around 2004, put ADCs, DACs and an FPGA on a USB~2.0 board with
plug-in RF daughterboards, at a price a university lab could afford. Ettus Research was acquired by
National Instruments in 2010, and the USRP family---including the B200 and B210 (introduced in 2013,
built around Analog Devices' AD9364 and AD9361 RFICs)---became the standard research SDR. In 2012 the
Linux developer Antti Palosaari found that the Realtek RTL2832U chip in cheap DVB-T television dongles
could be put into a mode that streamed raw 8-bit I/Q samples to the host; the \texttt{rtl-sdr} project of
the Osmocom community turned it into a receiver covering roughly 24 to 1766\,MHz for about \$20. Within a
year hundreds of thousands of hobbyists, students and security researchers had an SDR. Today
software-defined base stations run on commodity servers (OpenAirInterface, srsRAN), and the Open RAN
movement is built on the same idea.
\end{history}

\mdcpairany{ch07_usrp1_fpga}{The heart of the original USRP: an Altera Cyclone FPGA flanked by two
Analog Devices AD9862 mixed-signal front ends (ADCs and DACs) and a Cypress USB controller. Plug-in
daughterboards supplied the RF.}{ch07_rtlsdr}{A DVB-T television dongle opened up: the RTL2832U
demodulator chip (centre), an R820T tuner (left, near the antenna socket) and a 28.8\,MHz crystal. Put
into its test mode, it became the \$20 SDR that brought radio to the masses.}

\subsection{SDR architecture: RFIC, FPGA, host}

A modern SDR has three layers (Figure~\ref{fig:ch07_b200}), like a restaurant with a kitchen, a pass
and a dining room: the RFIC prepares the raw ingredients, the FPGA plates and times them, and the host
is where they are actually consumed.
\begin{description}
\item[The RF front end and RFIC.] An integrated zero-IF transceiver (AD9361/AD9364, AD9371, LMS7002M)
or, at the high end, RF data converters with discrete or semi-discrete front ends. It covers a wide
tuning range with programmable analog filters, gains and synthesisers, and it already does a
considerable amount of DSP: decimation, interpolation, fixed filtering, DC and I/Q tracking, AGC.
\item[The FPGA.] Further rate conversion (CIC and half-band filters, Chapter~\ref{ch:dsp}), digital
down- and up-conversion with an NCO for fine tuning, timestamping of every sample, packetisation, and,
on larger devices, user-programmable DSP for functions that are too fast for the host.
\item[The host.] A PC, embedded ARM processor or GPU running the driver (UHD for Ettus radios) and the
application---GNU Radio, MATLAB, a C++ program or Python with NumPy. All of the waveform-specific
processing in this book's labs lives here.
\end{description}

\begin{figure}[htbp]\centering
\begin{tikzpicture}[
  blk/.style={draw=navy,fill=navy!6,thick,minimum height=0.62cm,minimum width=1.05cm,align=center,font=\sffamily\tiny},
  fpga/.style={draw=accent,fill=accent!6,thick,minimum height=0.62cm,minimum width=1.05cm,align=center,font=\sffamily\tiny},
  arr/.style={-{Stealth[length=1.4mm]},thick,navy}]
% chip outlines
\draw[navy,thick,dashed,rounded corners] (0.9,-0.55) rectangle (8.25,2.75);
\node[font=\sffamily\scriptsize\bfseries,text=navy] at (4.6,2.5) {AD9364 RFIC};
\draw[accent,thick,dashed,rounded corners] (8.45,-0.55) rectangle (10.75,2.75);
\node[font=\sffamily\scriptsize\bfseries,text=accent] at (9.6,2.5) {Spartan-6 FPGA};
% RX row
\node[font=\sffamily\tiny,align=center] (rxp) at (0.0,1.5) {RX2 or\\TX/RX\\SMA};
\node[blk] (lna) at (1.6,1.5) {LNA\\(gain steps)};
\node[blk] (mix) at (2.85,1.5) {quadrature\\mixer};
\node[blk] (tia) at (4.1,1.5) {TIA +\\1st-order LPF};
\node[blk] (bbf) at (5.35,1.5) {BB LPF\\0.2--56 MHz};
\node[blk] (adc) at (6.6,1.5) {12-bit\\ADCs};
\node[blk] (dec) at (7.7,1.5) {HB + FIR\\decim.};
\node[fpga] (fr) at (9.6,1.5) {DC, CIC, HB\\decim., time-\\stamps};
\draw[arr] (rxp)--(lna); \draw[arr] (lna)--(mix); \draw[arr] (mix)--(tia); \draw[arr] (tia)--(bbf);
\draw[arr] (bbf)--(adc); \draw[arr] (adc)--(dec); \draw[arr] (dec)--(fr);
% TX row
\node[font=\sffamily\tiny,align=center] (txp) at (0.0,0.15) {TX/RX\\SMA};
\node[blk] (txo) at (1.6,0.15) {TX output\\atten.};
\node[blk] (upc) at (2.85,0.15) {quadrature\\modulator};
\node[blk] (tbf) at (4.1,0.15) {BB LPFs};
\node[blk] (dac) at (5.35,0.15) {12-bit\\DACs};
\node[blk] (int) at (6.6,0.15) {FIR + HB\\interp.};
\node[fpga] (ft) at (9.6,0.15) {CIC, HB\\interp.,\\timed bursts};
\draw[arr] (txo)--(txp); \draw[arr] (upc)--(txo); \draw[arr] (tbf)--(upc); \draw[arr] (dac)--(tbf);
\draw[arr] (int)--(dac); \draw[arr] (ft)--(int);
% synths
\node[blk,draw=practc,fill=practc!8] (syn) at (2.85,-1.25) {RX and TX\\frac-$N$ synth.};
\node[blk,draw=practc,fill=practc!8] (bbpll) at (5.35,-1.25) {BB PLL\\(converter clock)};
\node[blk,draw=practc,fill=practc!8] (tcxo) at (0.6,-1.25) {40 MHz TCXO\\(10 MHz/PPS,\\GPSDO opt.)};
\draw[arr,practc] (tcxo)--(syn); \draw[arr,practc] (syn)--(mix); \draw[arr,practc] (syn)--(upc);
\draw[arr,practc] (syn.east)--(bbpll.west); \draw[arr,practc] (bbpll)--(dac); \draw[arr,practc] (bbpll.east) -| (adc.south);
% host
\node[draw=navy,thick,fill=navy!12,minimum height=2.0cm,minimum width=1.4cm,align=center,font=\sffamily\tiny] (host) at (12.1,0.8) {USB 3.0\\$\downarrow$\\host PC:\\UHD driver,\\GNU Radio,\\Python};
\draw[arr] (fr.east) -- (fr.east -| host.west);
\draw[arr] (ft.east -| host.west) -- (ft.east);
\end{tikzpicture}
\caption{The signal chain of the USRP B200. The AD9364 RFIC does the analog work and the first stages
of rate conversion; the FPGA completes the decimation or interpolation, removes residual DC, timestamps
samples and packs them for the USB~3.0 link to the host, where UHD delivers them to the application. One
40\,MHz reference clocks everything. Specifications are approximate; see the Ettus and Analog Devices
documentation.}
\label{fig:ch07_b200}
\end{figure}

\subsection{The USRP B200 in detail}

If you have a B200 on your bench, everything in the first half of this chapter is inside it
(Figure~\ref{fig:ch07_b210}). Let us take it apart, block by block, with the vocabulary we now have.

\mdcpairany{ch07_b210}{A USRP B210 board, the two-channel sibling of the B200 (Ettus product photo).
SMA connectors for its RF ports line the edges, and the USB~3.0 cable that carries both power and
samples plugs in at the left.}{ch07_usrp_b100}{The USRP B100, the B200's USB~2.0 predecessor, in its
case with SMA cables on its two RF ports and the 10\,MHz reference and PPS inputs at the left: the same
arrangement of ports a B200 user meets today.}

\paragraph{The RFIC.} The AD9364 is the single-channel member of Analog Devices' AD9361 family. Its
receive path (top row of Figure~\ref{fig:ch07_b200}) is the zero-IF chain of
Section~\ref{sec:ch07:arch}: an LNA with gain steps, a passive quadrature mixer, a transimpedance
amplifier whose feedback network forms a first low-pass pole, a programmable third-order baseband
low-pass filter whose bandwidth can be set from about 200\,kHz to 56\,MHz, and 12-bit ADCs that run at a
multiple of the output rate. A digital chain of fixed half-band decimators (and an optional
decimate-by-3 stage) and a programmable FIR of up to 128 taps brings the rate down to at most
61.44\,MS/s while providing the sharp channel filtering that the gentle analog filters do not. The
transmit path mirrors it, with a programmable FIR and half-band interpolators ahead of 12-bit DACs,
baseband filters, a quadrature modulator and an output stage with about 90\,dB of attenuation range in
0.25\,dB steps. Two fractional-$N$ synthesisers with integrated VCOs provide the RX and TX LOs from 70\,MHz
to 6\,GHz, and a third PLL (the BB~PLL) generates the converter clock. On top of the signal path sit the
calibrations and tracking loops of Section~\ref{sec:ch07:iq}, an AGC with fast and slow modes, and an
internal state machine that sequences them.

\paragraph{The board.} Around the RFIC the B200 adds RF switches and baluns connecting the two SMA ports
(\texttt{TX/RX}, which can transmit or receive, and \texttt{RX2}, receive only), a 40\,MHz TCXO
(about $\pm2$\,ppm) with a PLL that can discipline it to an external 10\,MHz reference, a PPS input, an
optional on-board GPSDO, a Xilinx Spartan-6 FPGA, and a Cypress FX3 USB~3.0 controller. The board is bus
powered (an external supply is recommended for transmitting at full power) and Ettus quotes the maximum
output power at more than $+10$\,dBm and the receive noise figure at under 8\,dB.

\mdcfig{ch07_offset_tuning}{Offset tuning on a zero-IF SDR. Left: tuned exactly to the signal, the
residual DC spur (and the DC-removal notch) sit in its middle. Right: the RF LO is tuned 450\,kHz away
and the FPGA's NCO shifts the signal back digitally; the spur is now harmlessly outside the signal. A
UHD tune request with an LO offset does exactly this.}

\paragraph{Master clock rate and sample rates.} The \emph{master clock rate} (MCR) sets the AD9364's
data-interface rate, up to 61.44\,MHz for a single channel (30.72\,MHz per channel when a B210 runs two).
The host sample rate you request is the MCR divided by an integer decimation (or multiplied by an integer
interpolation) in the FPGA. UHD picks an MCR automatically, but choosing it yourself is often better.
Choosing an MCR that makes the FPGA's rate change even, or a power of two, gives the best filter
performance, because the CIC and half-band stages are designed for those ratios; choosing an MCR equal to
a multiple of the signal's natural rate (for example 30.72\,MHz for LTE and NR, whose sample rates are
multiples of 1.92\,MHz) avoids fractional resampling altogether; and a rate change of 1 (MCR equal to the
sample rate) bypasses the FPGA's filters entirely, leaving only the AD9364's.

\paragraph{The USB link.} Samples cross USB~3.0 as 16-bit I and 16-bit Q integers (\texttt{sc16}, four
bytes per complex sample), or as 12- or 8-bit samples to save bandwidth (Figure~\ref{fig:ch07_usb_rates}).
At 61.44\,MS/s \texttt{sc16} is 246\,MB/s, close to what a good USB~3.0 controller sustains in practice;
the usable rate depends on the host's controller and operating system, and Ettus specifies up to
56\,MHz of instantaneous bandwidth for the B200. If you need the full rate reliably, use a short cable, a
dedicated USB~3.0 port and the \texttt{sc8} or \texttt{sc12} wire format.

\mdcfig[0.5]{ch07_usb_rates}{USB payload against sample rate for the three wire formats. At the
B200's top rate, \texttt{sc16} needs about 246\,MB/s; \texttt{sc12} and \texttt{sc8} cut that by a
quarter and a half at the cost of resolution.}

\subsection{UHD and the host side}

UHD (the USRP Hardware Driver) is Ettus's open-source driver and API, in C++ with Python bindings. It
finds and configures the device, tunes it, sets gains, rates and antennas, and moves samples between the
radio and the application as \emph{streams}. With a B200 plugged in, the whole of
Section~\ref{sec:ch07:metrics} is one Python call away. A minimal capture looks like this:
\begin{lstlisting}
import uhd, numpy as np
usrp = uhd.usrp.MultiUSRP("type=b200")
x = usrp.recv_num_samps(1_000_000,   # number of samples
                        2.437e9,     # centre frequency (Hz)
                        10e6,        # sample rate (S/s)
                        [0],         # channel list
                        40)          # receive gain (dB)
x = x[0]                             # complex64 array, full scale = 1.0
\end{lstlisting}
Behind that call are the objects every real application uses: a \emph{tune request}, which splits a
requested frequency between the RF synthesiser and the FPGA's NCO (so that an LO offset can move the DC
spur out of the band); \emph{stream arguments}, which choose the host format (\texttt{fc32} complex float
or \texttt{sc16}) and the wire format; a \emph{receive streamer} that delivers buffers of samples with
\emph{metadata}---a timestamp for the first sample, and error codes; and \emph{timed commands}, which
schedule a tune, a gain change or a transmission at a future value of the radio's sample clock. Utility
programs (\texttt{uhd\_find\_devices}, \texttt{uhd\_usrp\_probe}, \texttt{uhd\_fft},
\texttt{rx\_samples\_to\_file}, \texttt{benchmark\_rate}) are the first things to run on a new setup.

\subsection{Streaming, buffering and time}

The link between radio and host is a stream of packets, and real-time streaming is a contract: the host
must consume receive samples as fast as the radio produces them and supply transmit samples before the
DAC needs them. The radio is a tap that never stops running; the host is a bucket-emptier who must
never fall behind (Figure~\ref{fig:ch07_overflow}). When it fails, UHD prints a single letter to the
console, and every SDR user learns to read them:
\begin{description}
\item[\texttt{O} (overflow)] the radio's receive buffer filled because the host did not read fast enough;
samples were \emph{lost}, and the metadata reports a discontinuity.
\item[\texttt{D} (dropped packet)] a sequence error on the receive path, usually a network or USB problem.
\item[\texttt{U} (underflow)] the transmit buffer ran dry; the DAC sent zeros, and the transmitted
waveform has a gap.
\item[\texttt{L} (late)] a timed command or transmit burst arrived after its scheduled time.
\item[\texttt{S} (sequence error)] transmit packets were lost or reordered on the way to the radio.
\end{description}

\mdcfig{ch07_overflow}{Anatomy of an overflow. The radio fills its receive buffer at a steady rate; the
host empties it in bursts. When the host stalls (orange), the buffer fills, every further sample is
thrown away and UHD prints \texttt{O}. Note that the buffer recovers only slowly afterwards: a host that
only just keeps up has no spare capacity to catch up.}

The causes are usually mundane: a Python loop doing too much work per buffer, a CPU frequency governor
throttling the processor, a USB controller shared with other devices, a laptop on battery power. The
remedies are to do less per sample in the hot path (vectorise with NumPy, move work to compiled blocks),
to enlarge the transport buffers (\texttt{num\_recv\_frames}, \texttt{recv\_frame\_size}), to use a lower
sample rate or a smaller wire format, and to test the host with \texttt{benchmark\_rate} before blaming
your own code.

\paragraph{Timestamps.} Every packet carries a timestamp from the radio's sample clock, which can be set
to match an external PPS. That enables timed transmission (send this burst at exactly time $t$), coherent
operation of several radios sharing a 10\,MHz reference and PPS, time-of-arrival measurements, and
tagging every capture with absolute time. It also means that ``when did this sample arrive?'' has an exact
answer, which a host clock can never give.

\paragraph{Latency.} The round trip from antenna to host and back through USB, the operating system and
the application's buffers is typically a few milliseconds. That is fine for most experiments, and even
for the 1\,ms subframes of an LTE base station implemented on a PC (srsRAN and OpenAirInterface do it with
careful real-time tuning), but far too much for protocols such as Wi-Fi, whose acknowledgement must start
16\,$\mu$s after the end of a received frame. Such protocols require the time-critical MAC functions to run
in the FPGA, close to the converters.

\mdcpairany{ch07_grc}{GNU Radio Companion with a QPSK synchronisation flowgraph: random source, pulse
shaping, channel model, then frequency-locked loop, polyphase clock sync and Costas loop, with Qt GUI
sinks to watch the signal at each stage---Chapter~\ref{ch:sync} as boxes and arrows.}{ch07_gr_qtgui}{A
running GNU Radio flowgraph's Qt GUI: time-domain I and Q, a constellation, and the transmitted and
received spectra. The same widgets appear when you run the book's hardware scripts.}

\subsection{GNU Radio}

GNU Radio models a radio as a \emph{flowgraph}: a directed graph of \emph{blocks} connected by
\emph{streams} of items (complex samples, floats, bytes). If UHD is the plumbing, GNU Radio is the
water-treatment plant: samples flow in at one end, through filters, mixers and demodulators, and out
as audio, packets or pictures at the other. Each block implements a \texttt{work} function
that consumes input items and produces output items; \emph{sync} blocks produce one output per input,
\emph{decimators} and \emph{interpolators} a fixed ratio, \emph{general} blocks any number. The runtime
gives each block its own thread and connects them with lock-free ring buffers, so a flowgraph is
naturally pipelined across CPU cores. Two further mechanisms carry information that does not fit in a
stream:
\begin{itemize}
\item \textbf{Stream tags} attach key--value metadata to a particular item in a stream: the UHD source
tags the first sample with \texttt{rx\_time} and \texttt{rx\_freq} and re-tags after an overflow or retune;
a burst transmitter uses \texttt{tx\_sob}/\texttt{tx\_eob} or a \texttt{packet\_len} tag to delimit
bursts; a synchroniser can tag the start of each detected frame.
\item \textbf{Messages} are asynchronous objects (polymorphic types, PMTs) passed between blocks outside
the streams, used for packets (a protocol data unit with its metadata), for control (retune the radio) and
for connecting the physical layer to MAC-layer logic.
\end{itemize}
Blocks are grouped into \emph{hierarchical blocks}, and new ones are written in C++ or Python as
\emph{out-of-tree} (OOT) modules. GNU Radio Companion (GRC) draws a flowgraph graphically and generates a
Python program from it (Figure~\ref{fig:ch07_grc}). The hardware scripts in this book's companion code
(Appendix~\ref{appB}) are written directly in Python in the same style, so that every block and parameter
is visible; Figure~\ref{fig:ch07_gr01} shows the one you will use first. GNU Radio 4.0, under development
in 2026, rebuilds the runtime around a new scheduler and a compile-time C++ block model, but the
flowgraph concept is unchanged.

\begin{figure}[htbp]\centering
\begin{tikzpicture}[node distance=0.5cm and 0.7cm,
  blk/.style={draw=navy,fill=navy!6,thick,rounded corners=2pt,minimum height=0.85cm,minimum width=2.0cm,align=center,font=\sffamily\scriptsize},
  hw/.style={draw=accent,fill=accent!6,thick,rounded corners=2pt,minimum height=0.85cm,minimum width=2.0cm,align=center,font=\sffamily\scriptsize},
  gui/.style={draw=practc,fill=practc!8,thick,rounded corners=2pt,minimum height=0.85cm,minimum width=2.0cm,align=center,font=\sffamily\scriptsize},
  arr/.style={-{Stealth[length=1.8mm]},thick,navy}]
\node[hw] (src) {USRP source\\(B200, \texttt{fc32})\\or \texttt{-{}-sim} tones};
\node[gui,right=1.2cm of src,yshift=1.3cm] (fft) {Qt frequency sink\\(spectrum)};
\node[gui,right=1.2cm of src,yshift=0.2cm] (wf) {Qt waterfall sink};
\node[blk,right=1.2cm of src,yshift=-0.9cm] (file) {head + file sink\\(\texttt{.cfile} + SigMF)};
\node[blk,right=1.2cm of src,yshift=-2.0cm] (mag) {$|x|^2$ +\\moving average};
\node[blk,right=of mag] (probe) {probe:\\power in dBFS};
\node[gui,left=0.9cm of src,align=center,minimum width=1.6cm] (sl) {sliders:\\centre freq.,\\RX gain};
\draw[arr] (src.east) -- ++(0.5,0) |- (fft.west);
\draw[arr] (src.east) -- ++(0.5,0) |- (wf.west);
\draw[arr] (src.east) -- ++(0.5,0) |- (file.west);
\draw[arr] (src.east) -- ++(0.5,0) |- (mag.west);
\draw[arr] (mag)--(probe);
\draw[arr,dashed,practc] (sl)--node[above,font=\tiny]{set\_freq}node[below,font=\tiny]{set\_gain}(src);
\end{tikzpicture}
\caption{The flowgraph of \lab{gnuradio/gr01\_spectrum\_iq\_capture.py}, the book's first hardware
script for the B200. One USRP source feeds a live spectrum and waterfall, an optional I/Q recorder that
writes a \texttt{.cfile} with a SigMF metadata file, and a power meter. Sliders retune the radio and set
its gain while it runs; \texttt{-{}-no-dc-corr} turns off UHD's DC-offset correction.}
\label{fig:ch07_gr01}
\end{figure}

\begin{tryit}[Your first B200 capture]
With a B200 and an antenna on the \texttt{RX2} port, run
\texttt{python gr01\_spectrum\_iq\_capture.py -{}-freq 100e6 -{}-rate 2e6} and you will see the FM
broadcast band in the spectrum and waterfall. Drag the \emph{RX gain} slider up until the spectrum
grows a forest of new signals---that is ADC clipping, and the pitfall of Section~\ref{sec:ch07:agc}.
Then rerun with \texttt{-{}-no-dc-corr} and look for the zero-IF DC spur at the centre. No hardware?
Add \texttt{-{}-sim}: the script substitutes two tones, a DC spur and an I/Q image so you can practise
the same diagnoses. Lab~1's \emph{A live waterfall} experiment does the same in pure Python.
\end{tryit}

\subsection{Calibration and frequency accuracy in practice}

An SDR delivers its best performance only after its imperfections are measured and corrected. On a B200:
\begin{itemize}
\item \textbf{DC offset and I/Q balance} are calibrated by the AD9364 at each tune and tracked in the
background; leave the automatic corrections on unless you are studying the impairments, and use offset
tuning when your signal has energy at DC.
\item \textbf{Gain} is not calibrated in absolute terms: a gain setting of 40\,dB does not mean a known
number of dBFS per dBm, and the relation varies by a few decibels across frequency and between units. For
power measurements, calibrate with a signal generator at the frequencies you care about.
\item \textbf{Frequency} is as good as the reference: about $\pm2$\,ppm from the TCXO, parts per billion
with an external 10\,MHz reference from a lab standard or with the GPSDO module once it has locked to
satellites. Receivers of real signals can of course estimate and correct the offset (Chapter~\ref{ch:sync});
the reference matters when the radio itself is the measuring instrument or when two radios must agree.
\end{itemize}

\mdcfig{ch07_nf_measure}{Measuring a B200's noise figure with nothing more than a terminator and a
signal generator (the numbers of the in-practice box, simulated). A $-80$\,dBm tone fixes the scale
between dBm and dBFS; the terminated floor, converted to dBm/Hz, sits 5\,dB above $kT_0$.}

\begin{inpractice}[Measuring a B200: noise figure, spurs and clipping]
Three measurements every SDR owner should make once.

\emph{Noise figure.} Terminate the receive port in 50\,$\Omega$, set a gain, and measure the noise power
spectral density in dBFS/Hz (average many FFTs and correct for the window's equivalent noise bandwidth).
Then inject a known tone from a signal generator, say $-80$\,dBm, and read its level in dBFS: the
difference is the gain in dBFS per dBm. Suppose the tone reads $-22$\,dBFS (a conversion of $+58$\,dB)
and the terminated noise floor reads $-111$\,dBFS/Hz: referred to the input that is $-169$\,dBm/Hz, so the
noise figure is $-169-(-174)=5$\,dB (Figure~\ref{fig:ch07_nf_measure}). Repeat across frequency and gain to
map the receiver; at low gain settings you will see the noise figure rise steeply as the ADC's own noise
takes over.

\emph{Spurs.} With the input terminated, look for lines: the DC spur (with correction off), the image of
any strong signal, and fixed spurs related to the master clock, the reference and the USB controller,
which stay at fixed absolute frequencies or fixed offsets from the LO as you retune. Knowing them prevents
you from ``discovering'' them in a capture.

\emph{Clipping.} Feed a strong signal and raise the gain until the I/Q samples approach $\pm1$; the
spectrum will grow a forest of intermodulation products and the amplitude histogram will pile up at the
rails. Back off by at least the signal's PAPR plus a few decibels. A good habit is to log the peak
magnitude of every capture.
\end{inpractice}

\subsection{The SDR landscape}
\label{sec:ch07:platforms}

Table~\ref{tab:ch07:platforms} lists representative platforms. At the low end, receive-only dongles and
the half-duplex HackRF (Figure~\ref{fig:ch07_hackrf}) serve hobbyists and teaching; the ADALM-Pluto,
with an AD9363 and a small Zynq, is built for the classroom. The B200/B210 and bladeRF occupy the
research middle ground, with full-duplex AD9361-family RFICs. At the top, radios built on the Zynq
UltraScale+ RFSoC (such as the USRP X410) sample directly at RF with hundreds of megahertz of bandwidth
per channel, and platforms such as Deepwave's AIR-T put an embedded GPU next to the RFIC for
machine-learning signal processing at the edge. The specifications are approximate and evolve with each
hardware revision.

\begin{table}[htbp]\centering\small
\caption{Representative SDR platforms (approximate specifications).}
\label{tab:ch07:platforms}
\setlength{\tabcolsep}{4pt}\footnotesize
\begin{tabular}{lllll}
\toprule
Device & RF chip & Frequency range & Bandwidth & Notes \\
\midrule
RTL-SDR dongle & RTL2832U + tuner & $\approx$24\,MHz--1.7\,GHz & $\approx$2.4\,MHz & 8-bit, receive only \\
HackRF One & MAX2837 + mixer & 1\,MHz--6\,GHz & 20\,MHz & 8-bit, half-duplex \\
ADALM-Pluto & AD9363 & 325\,MHz--3.8\,GHz & 20\,MHz & 1$\times$1, Zynq-7010, teaching \\
USRP B200 / B210 & AD9364 / AD9361 & 70\,MHz--6\,GHz & 56\,MHz & 1$\times$1 / 2$\times$2, USB 3.0 \\
bladeRF 2.0 micro & AD9361 & 47\,MHz--6\,GHz & 56\,MHz & 2$\times$2, USB 3.0 \\
LimeSDR & LMS7002M & 0.1\,MHz--3.8\,GHz & 61.44\,MHz & 2$\times$2, open hardware \\
AIR-T & AD9371 & 300\,MHz--6\,GHz & $\approx$100\,MHz & 2$\times$2, embedded GPU \\
USRP X410 & RFSoC & 1\,MHz--7.2\,GHz & 400\,MHz & 4$\times$4, direct sampling \\
\bottomrule
\end{tabular}
\end{table}

\mdcphoto[0.62]{ch07_esa_chamber}{A prototype antenna under test in ESA's Compact Antenna Test
Facility, a shielded anechoic chamber whose foam pyramids absorb reflections. Professionals test
transmitters where their signals cannot escape; for an SDR experimenter, the equivalent is a coaxial
cable and a 40\,dB attenuator between TX and RX.}

\begin{inpractice}[SDR in industry]
\emph{Base stations.} Every modern base station is software-defined in its baseband, and increasingly in
its radio: an O-RAN radio unit using the 7.2x functional split performs the FFT, cyclic prefix, part of the
beamforming, digital up- and down-conversion, CFR and DPD in an FPGA or ASIC next to the RF, while the
distributed unit runs the rest of the physical layer on servers, sometimes with accelerator cards
(Chapter~\ref{ch:cellular}). \emph{Defence.} Tactical radios in NATO forces are SDRs running portable
waveforms under the SCA and its successors, and electronic-warfare systems rely on wideband
direct-sampling receivers. \emph{Spectrum monitoring.} National regulators and operators use networks of
wideband SDR receivers to find interference and unlicensed transmitters, and the same hardware does
direction finding by time difference of arrival. \emph{Amateur radio and broadcasting.} Direct-sampling SDR
transceivers have taken over the high end of the amateur market, and broadcast exciters for DAB, DVB-T2 and
ATSC~3.0 are software-defined. \emph{Space and science.} Satellite ground stations, including volunteer
networks such as SatNOGS, CubeSat radios, GNSS research receivers (Chapter~\ref{ch:spreadspectrum}) and
radio-astronomy back ends are SDRs. \emph{Research and test.} Open-source stacks (srsRAN,
OpenAirInterface) on USRPs are the standard way to prototype 5G and 6G ideas, and the same radios are
used to test the security of everything from car key fobs to aircraft transponders.
\end{inpractice}

\begin{pitfall}[Transmitting legally and safely]
A B200 can transmit anywhere from 70\,MHz to 6\,GHz, including frequencies used by aviation, GNSS,
public-safety and cellular systems; transmitting there without authorisation is illegal almost
everywhere and can endanger lives. Over the air, transmit only where you are licensed (an amateur licence
covers specific bands and modes) or in licence-exempt ISM bands within the applicable power and duty-cycle
rules (in the US, FCC Part~15; in Europe, the relevant ETSI harmonised standards). For everything else, use
a cable. The B200 can deliver more than $+10$\,dBm, and its receiver is rated for a maximum input of about
$-15$\,dBm: never connect TX to RX without at least 30\,dB of attenuation (40 to 60\,dB is safer and leaves
the receiver in its linear range), start with the lowest TX gain, and remember that the harmonics and
spurious emissions of an unfiltered SDR transmitter may land in somebody else's band even when the
fundamental is legal.
\end{pitfall}

%======================================================================
\mdcphoto[0.7]{ch07_hackrf}{The HackRF One, an open-hardware SDR. The RF section is the shielded area at
lower left: a MAX2837 2.4\,GHz transceiver chip plus an extra mixer stage that extends it from 1\,MHz to
6\,GHz. Above it sits a small Xilinx CPLD and, at the right, the NXP microcontroller that streams 8-bit
samples over USB. Every architecture in this chapter can be spotted on boards like this one.}

\section*{Summary}
\addcontentsline{toc}{section}{Summary}
%======================================================================

\begin{itemize}
\item Receivers come in four families: the superheterodyne (best selectivity, needs image and IF filters),
zero-IF (integrable, but with DC offset, LO leakage, flicker noise, IM2 and I/Q imbalance), low-IF (away
from DC, but the image is a neighbour) and direct RF sampling (all-digital, power- and ADC-limited).
Hartley and Weaver structures reject the image by phasing.
\item The front end's filters, duplexer and LNA set most of the noise figure; loss before the LNA counts
in full. Switching mixers respond at odd LO harmonics; the Gilbert cell trades noise and linearity for gain.
\item Sensitivity is $-174+10\log_{10}B+\text{NF}+\text{SNR}_{\min}$\,dBm. Noise wants gain early, linearity
late; IIP3 cascades as the sum of $G_{\text{before}}/\text{IIP3}_k$; SFDR $=\tfrac23(\text{IIP3}-P_{\text{noise}})$.
Blockers desensitise by compression, AGC, IM2 and reciprocal mixing; a level diagram and analog filtering
before the ADC make the dynamic range close.
\item References range from $\pm20$\,ppm XOs to GPSDOs; fractional-$N$ PLLs with sigma-delta dividers give
fine steps with wide loops. Leeson's model explains the $1/f^3$, $1/f^2$ and flat regions of phase noise;
integrated phase noise limits QAM order and OFDM subcarrier spacing; reciprocal mixing,
$P_b+\mathcal L(\Delta f)+10\log_{10}B$, couples blockers into the channel.
\item I/Q imbalance creates a conjugate image with $\text{IRR}\approx4/(\epsilon^2+\phi^2)$; $0.1$\,dB and
$1^\circ$ give about 40\,dB. Blind second-order-statistics estimators and loopback calibration correct it,
and offset tuning sidesteps DC.
\item PA efficiency collapses with back-off; OFDM's PAPR makes linear PAs inefficient. ACLR and EVM bound
the distortion. CFR, DPD with memory polynomials and indirect learning, Doherty load modulation and envelope
tracking together make linear modulation affordable.
\item An SDR is RFIC + FPGA + host. The B200 couples an AD9364 zero-IF transceiver to a Spartan-6 FPGA and
USB~3.0; master clock rate, wire format, buffering and timestamps determine whether samples arrive intact.
GNU Radio builds radios as flowgraphs of blocks with stream tags and messages. Measure your radio's noise
figure, spurs and clipping point before trusting it, and transmit only where you may.
\end{itemize}

\begin{labbox}[The imperfect transceiver]
{\raggedright Lab files: \lab{moderncomms-labs/labs/lab36\_rf\_transceiver.py},
\lab{lab01\_baseband.py}, \lab{lab15\_superhet\_receiver.py}, \lab{lab28\_ofdm\_system.py} and the
hardware script \lab{moderncomms-labs/gnuradio/gr01\_spectrum\_iq\_capture.py}\par}
\smallskip
\noindent\textbf{Lab 36} (run \lab{python lab36\_rf\_transceiver.py}) is the RF engineer's bench for this
chapter, built on \texttt{commlib.rf}, the same models as the figures:
\begin{itemize}
\item \emph{Two-tone test and IIP3}: cubic and soft-limiter amplifiers, IM3 slope and intercept, P1dB
(Figure~\ref{fig:ch07_two_tone}).
\item \emph{Receiver line-up builder}: cascaded gain, NF, IIP3 and SFDR stage by stage, as in the worked
example (Figures~\ref{fig:ch07_whisper_chain} and~\ref{fig:ch07_cascade_tradeoff}).
\item \emph{Phase noise and reciprocal mixing}: an $\mathcal L(f)$ mask, a blocker and the rise of the
noise floor (Figure~\ref{fig:ch07_reciprocal}).
\item \emph{I/Q imbalance and the image}: skewed constellations, mirror neighbours and blind correction
(Figures~\ref{fig:ch07_irr} and~\ref{fig:ch07_iq_correction}).
\item \emph{Power amplifier: ACLR and EVM}: regrowth and EVM against back-off with 3GPP limits
(Figures~\ref{fig:ch07_pa_regrowth} and~\ref{fig:ch07_aclr_backoff}).
\item \emph{PA classes, Doherty and ET}: average efficiency with real signal statistics
(Figures~\ref{fig:ch07_pa_efficiency} and~\ref{fig:ch07_doherty}).
\item \emph{Crest-factor reduction}: clip-and-filter CCDF and spectrum (Figure~\ref{fig:ch07_cfr}).
\item \emph{Digital predistortion}: indirect learning one iteration per click, with and without PA memory
(Figure~\ref{fig:ch07_dpd}).
\end{itemize}
\textbf{Lab 1} walks down the B200/AD9364 receive chain: the IQ helix, real versus complex sampling,
up- and down-conversion, a digital down-converter, \emph{Zero-IF fingerprints} (DC offset, I/Q imbalance and
the IRR of~\eqref{eq:ch07:irr}, with blind correction by~\eqref{eq:ch07:iqcorr}), oscillator phase noise,
ADC bits and processing gain, sensitivity and range, and a live waterfall. \textbf{Lab 15} builds a
superheterodyne stage by stage (image problem, IF selectivity, AGC, double conversion, intermodulation and
SFDR) and compares zero-IF with low-IF. \textbf{Lab 28}, Section~4, clips, filters and amplifies an OFDM
signal through a PA model and measures EVM and ACLR. With hardware, \lab{gr01\_spectrum\_iq\_capture.py}
is a receive-only B200 spectrum, waterfall and I/Q recorder (Figure~\ref{fig:ch07_gr01});
\texttt{-{}-no-dc-corr} reveals the zero-IF DC spur and \texttt{-{}-sim} runs without hardware.
\end{labbox}

\problems
\begin{enumerate}[label=\thechapter.\arabic*]
\item A superheterodyne receiver for the 88--108\,MHz FM band uses a 10.7\,MHz IF with high-side LO.
What range must the LO tune over, and where do the images of stations at 88 and 108\,MHz lie? Why is the
image not in the FM band? What image rejection does a single tuned circuit with loaded $Q=30$ at
the signal frequency provide for a station at 98\,MHz?
\item A receiver has an LNA (NF 1.5\,dB, gain 18\,dB, IIP3 $+5$\,dBm), a mixer (NF 9\,dB, gain 8\,dB,
IIP3 $+10$\,dBm) and an IF amplifier (NF 6\,dB, gain 30\,dB, IIP3 $+15$\,dBm). Compute the cascaded NF
and IIP3 and the SFDR in a 1\,MHz bandwidth. Which stage limits each?
\item Derive the SFDR formula~\eqref{eq:ch07:sfdr} and the cascaded IIP3 formula~\eqref{eq:ch07:iip3cascade}.
Under what assumption about the phases of the IM3 products is the latter a worst case?
\item Two interferers at 1001 and 1002\,MHz, each $-35$\,dBm, reach a receiver tuned to 1000\,MHz with
IIP3 $-5$\,dBm. What is the IM3 level on the wanted channel? By how much must the interferers fall for
the IM3 product to drop 20\,dB?
\item For $y=a_1x+a_3x^3$, show that the 1\,dB compression point is at $A^2=0.145|a_1/a_3|$ and the
intercept at $A^2=\tfrac43|a_1/a_3|$, so that IIP3 is 9.6\,dB above P1dB.
\item Explain why a zero-IF receiver is sensitive to IIP2 but a superheterodyne largely is not. An FDD
phone's transmitter leaks $-27$\,dBm of a 20\,MHz-wide OFDM signal into its receiver. What IIP2 is needed
for the IM2 product to be 10\,dB below a $-105$\,dBm noise floor? (Treat the IM2 product as noise-like.)
\item Derive the I/Q imbalance model~\eqref{eq:ch07:iqmodel} and the IRR~\eqref{eq:ch07:irr}. What phase
error alone gives 30\,dB? What gain error alone? A low-IF receiver must handle a 35\,dB stronger image
channel with 12\,dB SIR: what IRR does it need, and is it achievable without calibration?
\item A 3.5\,GHz receiver's LO has $\mathcal{L}=-100$\,dBc/Hz at 10\,kHz, falling 20\,dB/decade to a floor of
$-150$\,dBc/Hz. Compute the RMS phase error integrated from 10\,kHz to 10\,MHz and the resulting EVM
contribution. Is it adequate for 256-QAM? What changes if the carrier is 28\,GHz with the same synthesiser
design?
\item A GSM receiver (200\,kHz channel, NF 6\,dB) must tolerate a $-23$\,dBm blocker at 3\,MHz with no more
than 3\,dB of desensitisation. What LO phase noise at 3\,MHz offset is required? Repeat for a $-43$\,dBm
blocker at 600\,kHz.
\item An OFDM PA delivers 20\,W average with 8\,dB of back-off and 25\% efficiency. Estimate the DC
power and annual energy saved if DPD allows 5\,dB of back-off at 38\% efficiency, and if a Doherty
architecture then raises the efficiency at that back-off to 55\%.
\item Using~\eqref{eq:ch07:doherty}, find the efficiency of an ideal symmetric Doherty at its minimum
between 0 and 6\,dB back-off, and show that it occurs at $v=2/3$. For a Rayleigh-envelope signal clipped
at a PAPR of 8\,dB, compute the average efficiency of class B and of the symmetric Doherty numerically.
\item Choose a B200 master clock rate for (a) a 1\,MS/s experiment, (b) a 20\,MHz LTE carrier sampled at
30.72\,MS/s, (c) a 25\,MS/s wideband capture. Justify each choice, and compute the USB throughput for
each with \texttt{sc16} samples.
\item \emph{(Simulation)} Implement a memory-polynomial DPD with indirect learning for a PA model with
memory (a Wiener--Hammerstein model), starting from the chapter's figure script (\lab{ch07\_figs.py}). Plot ACLR
against $K$ and $M$ and against input back-off, and compare with a memoryless DPD. At what back-off does
DPD stop helping, and why?
\item \emph{(Simulation)} Generate a 16-QAM signal at DC and a 20\,dB stronger carrier at $+0.3f_s$,
apply 1\,dB and $5^\circ$ of I/Q imbalance, and implement (a) the blind estimator~\eqref{eq:ch07:iqcorr}
and (b) an adaptive circularity-restoring corrector $z=y+wy^*$. Plot the residual image against the
number of samples used, and show where each fails.
\item \emph{(Hardware)} With a B200, a 50\,$\Omega$ terminator and a signal generator (or a second
SDR through 60\,dB of attenuation), measure the receiver's noise figure at 915\,MHz and 2.4\,GHz for gain
settings of 20, 40 and 60\,dB, following the in-practice box. Plot NF against gain and explain the shape.
\item \emph{(Hardware)} Using \lab{gr01\_spectrum\_iq\_capture.py} with and without
\texttt{-{}-no-dc-corr}, capture a strong FM broadcast station tuned (a) exactly to its carrier and (b)
300\,kHz away. Measure the DC spur level and the station's image, compute the IRR, and estimate the gain
and phase imbalance with Lab~1's \texttt{iq\_correct}.
\end{enumerate}

\furtherreading
\begin{itemize}
\item B.~Razavi, \emph{RF Microelectronics}, 2nd ed., Prentice Hall, 2012. The standard text on
transceiver architectures and circuits: noise, nonlinearity, receiver and transmitter architectures, LNAs,
mixers, oscillators, PLLs and PAs, with a system designer's eye throughout.
\item D.~M.~Pozar, \emph{Microwave and RF Design of Wireless Systems}, Wiley, 2001. A clear treatment of
the system-level quantities of this chapter---noise figure, intercepts, dynamic range, filters, mixers,
oscillators and link budgets---for readers coming from microwave engineering.
\item S.~C.~Cripps, \emph{RF Power Amplifiers for Wireless Communications}, 2nd ed., Artech House, 2006.
The practitioner's book on PA classes, load-line design, Doherty, envelope tracking and linearisation.
\item A.~A.~Abidi, ``Direct-conversion radio transceivers for digital communications,'' \emph{IEEE Journal
of Solid-State Circuits}, vol.~30, no.~12, 1995. The classic analysis of zero-IF impairments.
\item D.~B.~Leeson, ``A simple model of feedback oscillator noise spectrum,'' \emph{Proceedings of the
IEEE}, vol.~54, no.~2, 1966. Two pages that every oscillator designer has read.
\item J.~Mitola, ``The software radio architecture,'' \emph{IEEE Communications Magazine}, vol.~33,
no.~5, 1995. The paper that defined the field's ambitions.
\item J.~H.~Reed, \emph{Software Radio: A Modern Approach to Radio Engineering}, Prentice Hall, 2002. A
broad survey of SDR architecture, data converters, front ends and smart antennas from the era when SDR
became practical.
\item D.~R.~Morgan, Z.~Ma, J.~Kim, M.~G.~Zierdt and J.~Pastalan, ``A generalized memory polynomial model
for digital predistortion of RF power amplifiers,'' \emph{IEEE Transactions on Signal Processing},
vol.~54, no.~10, 2006. The memory-polynomial family and its generalisation, with measured results.
\item W.~H.~Doherty, ``A new high efficiency power amplifier for modulated waves,'' \emph{Proceedings of
the IRE}, vol.~24, no.~9, 1936. Still readable, and the origin of today's base-station PA.
\item Ettus Research, \emph{USRP B200/B210 data sheet} and the \emph{UHD manual} (online). Specifications,
clocking, master-clock-rate rules, stream arguments and the error codes of
Section~\ref{sec:ch07:sdr}.
\item Analog Devices, \emph{AD9361/AD9364 Reference Manual (UG-570)} and the AD9364 data sheet. The
internal signal chain, gain tables, filters, calibrations and AGC of the B200's RFIC, in exhaustive detail.
\item 3GPP TS~38.101-1 and TS~38.104, \emph{NR User Equipment and Base Station radio transmission and
reception}. The source of the sensitivity, blocking, ACLR, EVM and frequency-error requirements that drive
transceiver design.
\end{itemize}
